\documentclass[12pt,a4paper,openany]{book}
% English edition preamble — pdfLaTeX + utf8 + English babel
\usepackage[T1]{fontenc}
\usepackage[utf8]{inputenc}
\usepackage[british]{babel}

\usepackage{geometry}
\geometry{a4paper, top=2.5cm, bottom=2.5cm, left=3cm, right=2cm}

\usepackage{microtype}
\usepackage{indentfirst}
\usepackage{setspace}
\usepackage{parskip}
\emergencystretch=2em
\tolerance=1500
\setlength{\parindent}{1.25cm}
\setlength{\parskip}{0pt}

\usepackage{hyperref}
\hypersetup{
  colorlinks=true,
  linkcolor=black,
  urlcolor=blue,
  pdftitle={What Is Communism?},
  pdfauthor={Engels Group},
  pdfsubject={Theory of communism},
  pdfkeywords={communism, Marxism, Marx, Engels, Lenin},
  pdfencoding=auto,
  unicode=true,
}

\usepackage{etoolbox}

\usepackage{titlesec}
\titleformat{\chapter}[block]{\centering\large\bfseries}{}{0pt}{}
\titleformat{\section}[block]{\centering\normalsize\bfseries}{}{0pt}{}
\titlespacing*{\chapter}{0pt}{24pt}{12pt}
\titlespacing*{\section}{0pt}{18pt}{9pt}

\usepackage{graphicx}
\graphicspath{{./images/}{../extracted/images/}{./}{../}}

\usepackage{booktabs}
\usepackage{longtable}

\usepackage{xcolor}
\usepackage{fancyhdr}
\pagestyle{plain}

\usepackage{tocbibind}
\usepackage{tocloft}



\pretocmd{\section}{\phantomsection}{}{}

% No chapter/section numbering for this unnumbered work
\setcounter{secnumdepth}{0}

\begin{document}
\ifSubfilesClassLoaded{\setcounter{page}{4}}{}
\section{Inventing the “Bicycle Spectacles”}
\label{sec:01_inventing_spectacles}

What is communism? Where did it come from, where is it calling us, and whom is it calling? As History has shown, all these simple questions demand answers not once and for all, but every time they arise. As the screenwriter Grigori Gorin put it through the mouth of “The Very Same Munchhausen”, “love is a theorem that has to be proved every day”. It is the same for us, except that communism will be a theorem (a theory, even!) far more serious and far more global than love. And first we must learn to write down correctly, under the “given” of this theorem, everything that is genuinely present rather than merely imagined.

Objections to the very word “communism” are dictated by the historical regression our society has undergone. The vulgarisation of this concept after the collapse of the USSR, and even its transformation into a term of abuse, came about because of the triumph of social forces openly hostile to communism.

Even before the Soviet Union fell, the “blurring” of the concept had already begun, and so some theorists spoke of the need to draw distinctions within Soviet economics and ideology, to draw a dividing line so as not to mistake the inertial movement of capitalism for a communist tendency. Afterwards—because society had been plunged into the restoration of capitalism at a time when the tradition of communist theoretical thought had been all but lost, and the question was one of preserving at least something of it—the problems of communism and of the movement towards it were not specifically studied. All eyes were turned to flourishing capitalism. Both its supporters and its opponents spoke of it as something real, whereas communism and the movement towards it were seen as things of the past, or, in the case of its supporters, as a future so distant that one would first have to live long enough to see it, and so the important thing was to focus on surviving. For decades, communism as a theoretical problem ceased to exist.

Mass culture, far removed from theoretical thinking, has always actively reproduced (and still reproduces) the consciousness of the ruling class. This class fears and hates communism, clothing its fear and hatred in the form of various social myths on which a second generation is now growing up in the post-Soviet countries. Yet mythologised representations of communism have nothing in common with the concept of communism. Here we find the first crucial point of “hostility” (which is not actually hostility at all). Citizens of countries that slid into capitalism not because of foreign intervention but “under their own steam” simply no longer know what genuine, that is, scientific, communism is, or what the essence of the idea of communism is. And yet the problem of communism is once again confronting them as a practical one. And the integral theory whose trampling and neglect played no small part in the backslide into capitalism in the USSR and the other countries of the socialist camp needs to be revived and explained “in plain terms” to the people whose genuine interests it expresses. This applies not only to politically indifferent working proletarians, but also to those who consider themselves communists, or at least leftists—that is, to those who have politically defined themselves as sympathetic to communism. The latter perceive communism as some sort of “good” social order, interpreting this “good” rather at the level of sensory representations, often “by contradiction”—as the absence of the abominations and horrors of modern capitalist society. But such supporters cannot give a positive answer to the question “What is communism?” even to themselves.

This state of affairs is precisely what prompts us to produce this pamphlet now, as a pair of “bicycle spectacles” that allow one to see communism again and understand what it is. As for the opponents of communism, we must at the very least force them to raise not foolish but intelligent objections to communism—now scientific communism.

In the centenary year of October, interpretations of the role and conduct of the Bolsheviks in 1917 were so diverse in the public sphere that opinions and epithets predictably polarised: from noble heroes who “saved the state” to bandits who hated and destroyed that state (and statehood itself). But amid this polarisation the main point slipped away: the Bolsheviks, and of course Lenin, were thinkers, and in their actions they were guided not only by desires, preferences and so on, but also by a theory that was the culmination of the entire preceding development of human thought. The enemies of socialism and communism like to call October and the subsequent socialist construction an experiment performed by the Bolsheviks on the peoples… Well, these enemies are far closer to the truth than friends who are far removed from theory.

The Great October Revolution was precisely an experiment, provided, of course, that one does not understand the word foolishly, as some kind of external manipulation. Given how unprecedented it was, it could not have been anything else. It was an experiment conducted in keeping with every scientific tradition, for there is no experiment without a hypothesis. And there was such a hypothesis, developed by Lenin and his supporters and grounded in the legacy of Marx and Engels: the World Revolution. Having begun in Russia, the proletarian revolution was meant to stir up the industrially developed countries of Europe and set in motion their transformation into the World Republic of Soviets. By no means every scientific hypothesis is confirmed by experiment in its original formulation (even though the idea that proletarian revolutions are relevant remains broadly true today), yet even this experimental movement along a path charted by an even partly unconfirmed hypothesis moves science (socialism) forward.

The proposed metaphor simply brings us back to a different way of seeing the revolution, War Communism, the NEP, industrialisation and so on. It calls on us to view modern society and the trends in its development through the prism of communism as the most important science within the whole system of the social sciences that are needed in order to act rationally in today’s world. It is needed so that the revolutionary class can overcome its submission to constant exploitation, overcome within itself the oppressed class of capitalist society and take up its own class position, so that it can comprehend social relations and rationally change them towards the abolition of all oppression and of the exploitation of man by man, towards communism.

That is why the “bicycle spectacles” of the theory of communism, mocked by Mayakovsky in a different context, will only help us if we have the patience to “reinvent” them (to reproduce them): to travel Marx’s road to the Commune and the International, Lenin’s theoretical road to October, and then the entire Soviet road, without changing our “means of transport”—revolutionary theory.

To reinvent means to make the legacy of the science of communism our own under the present-day historical conditions we have inherited.
\ifSubfilesClassLoaded{\setcounter{page}{7}}{}
\section{Preface}
\label{sec:02_preface}

In his Critique of the Gotha Programme, Marx wrote: “What we have to deal with here is a communist society, not as it has developed on its own foundations, but, on the contrary, just as it emerges from capitalist society; which is thus in every respect, economically, morally and intellectually, still stamped with the birthmarks of the old society from whose womb it emerges”, and “In a higher phase of communist society, after the enslaving subordination of the individual to the division of labour, and therewith also the antithesis between mental and physical labour, has vanished; after labour has become not only a means of life but life’s prime want; after the productive forces have also increased with the all-round development of the individual, and all the springs of social wealth flow more abundantly—only then can the narrow horizon of bourgeois right be crossed in its entirety and society inscribe on its banners: From each according to his ability, to each according to his needs!”

The idea of private property and the rule of capital is opposed by the idea of communism. In reality, the only thing that can be set against a society in which capital rules is a communist society. But what is communism?

In this pamphlet we are concerned with what Marx called the higher phase of communism, that is, a communist society that has developed on its own foundations, or communism as such. The lower phase of communism, or the period of transition from capitalism to communism, which later came historically to be known as “socialism”, will be examined later. For now, the question is what exactly this transition is supposed to lead to, or what communists want to obtain as a result of the revolutionary transformation of social life.

What interests us here are the essential determinations that make a society specifically communist and distinguish it from all the modes of production and ways of conducting social life that preceded communism. We will therefore set aside any description of technology under communism. It may take various forms (and will, of course, change rapidly). For example, capitalism in the age of the steam engine and capitalism with modern digital technology differ enormously from one another. But what makes capitalism capitalism—a society of developed commodity production in which labour power, among other things, is a commodity, and capital is the dominant social relation—is characteristic, necessary and invariable for capitalism in general, in the 19th century and the 21st alike. An understanding of these necessary moments in their necessary interconnection is what constitutes the concept of capitalism.

In exactly the same way, the concept of communism, like that of any other object, includes only what is necessary and sufficient, in its necessary transitions into one another, for a society to be communist.

The difficulty lies in the fact that we are dealing with an object that has not yet become a reality. History has not yet known (at least on our planet) a communist society that has developed on its own foundations (primitive communism is irrelevant in this case).

However, communism long ago became a science. The entire course of the historical development of thought gave rise to the idea of communism, while the course of social life, at the very least, constantly demonstrates the limit of capitalism and the need to move to a new foundation of social development. Capitalist crises of overproduction are accompanied by the colossal destruction not only of the commodities produced or of individual enterprises, but often of entire once-thriving megacities.

If we consider social production as a mode of interaction between humans and nature, it becomes obvious that it is irrational and is inflicting harm on our planet on such a scale that the damage may become irreparable. The production and overproduction of that same packaging, which is excessive and pointless from the standpoint of a rational interaction with nature but necessary for the commodity to be sold, does enormous harm. The production of commodities themselves, which neither enrich nor develop people, and often even cripple them, yet consume vast quantities of resources and pollute the soil, water and air—this is the reality of capitalism. There is no well-thought-out waste recycling planned on the scale of society as a whole, nor any rational distribution and use of all natural resources and productive capacities. All of this creates the threat of an irreversible ecological catastrophe.

Alongside this man-made destruction of the natural world around us comes the degradation of humanity itself. This is not only the absolute poverty into which a substantial portion of the world's working people have been plunged: people go hungry, live in unsanitary conditions, and have no access to education or healthcare. It is also a matter of relative impoverishment, whereby new needs are produced whose satisfaction becomes a necessity even just to reproduce oneself as a seller of one's own labour power, yet the means to satisfy them are insufficient. Engels already noted this tendency, but today, when capitalism cannot survive without purposefully producing needs, it is intensifying many times over and assuming alarming proportions, especially in the periphery and semi-periphery of the capitalist world, where, on a societal scale, labour power is not fully restored. This undermines the very foundations of capitalist production.

But that is not all. In order to reproduce itself on an expanded scale, modern capital revives and subordinates to itself pre-capitalist forms of the exploitation of man by man. Slavery and serfdom exist in the modern world insofar as they prove profitable and contribute to the self-expansion of value. For example, according to ILO data published in 2014, profits from the use of slave labour worldwide exceed 150.2 billion dollars (as you can see, the profit from slavery is calculated, and can be calculated, in purely capitalist terms rather than in kind, as it was in the era of slaveholding).

Capital constantly demands room for self-expansion, and to achieve this it must push other capital out of existing markets and redistribute resources in its own favour. This generates conflicts not only within states but, above all, between them, which gives rise to local military actions and leads towards yet another world war. And when we have a tendency towards world war (its preparation, its unfolding, its individual flare-ups) combined with nuclear weapons, what accidents are not possible!… At the same time, war itself gives a powerful impetus to the development of a broad, mass, spontaneous anti-capitalist movement. How consistent this line will be, and whether it will lead to a genuine way beyond capitalism, is hard to say for now. In general, no one today can guarantee that humanity will outlive capitalism. The only thing this testifies to is the limit of development of a society in which capital dominates.

That capitalism will come to an end is beyond doubt. The only question is whether it will end together with humanity, or whether humanity will learn to live differently. And how exactly can we live differently, so that humanity as a whole and all its diverse components develop not by hindering but by helping one another? The scientific theory of communism provides the answer to this question, which is by no means outdated but, on the contrary, burningly relevant for humanity.

It is the main points of this theory that we will try to set out here. It was developed by Marx and Engels. Later Marxist thinkers also made their contribution to the development of the understanding of communism. Among them, Lenin's theoretical work deserves particular attention.

The meanings of the terms “communism” and “socialism” did not become established at once. For a long time they were used as synonyms, and communism was sometimes even understood as preceding socialism. Therefore, one should turn equally both to those works of Marx and Engels in which the term “communism” is used and to those that speak of “socialism”, provided it is clear from the substance of the text that what is being discussed is communist society, the communist movement, or the theory of communism. As for Lenin and other Marxists of the 20th century, in their works, as in Marx’s later works, the terminology had already become more or less settled. In his article “A Great Beginning”, Lenin specifically reminds readers that “the scientific distinction between socialism and communism is only that the first term means the first stage of the new society arising out of capitalism, while the second means the higher, further stage of it”. Wherever we use these words without references or special explanations, they should be understood according to contemporary usage.

As for the argument against communism that the various ideological servants of today’s capitalist ruling class keep repeating in every possible way—“communism is a utopia; people have already tried to build communism, and nothing (or nothing good) came of it”—our answer is as follows. Yes, it has not yet been possible to build a communist society, but to say that nothing good came of the attempt is simply a lie. In their day, the successes of the socialist Soviet Union could not be denied even by the most zealous anti-communists, even in the USA. Those who remember the recent history of the 20th century know very well that society’s movement towards communism enabled all the peoples living in the USSR to rise to heights never before seen under feudalism or capitalism, and to do so in the shortest possible time. This is borne out both by dry figures and calculations and by achievements in culture that remain unsurpassed to this day.

If we compare the first decades\footnote{This was despite the fact that the movement began in a war-ravaged, largely agrarian country surrounded by hostile powers; it was a movement along an untrodden path and therefore proceeded amid great problems, difficulties and retreats, and even tragedies, which we communists have no right to forget. Moreover, this movement took place in a society with rather strong capitalist tendencies, which cannot be abolished by decree or simply stepped over without doing away with commodity production.} in the life of a society moving towards communism with the same span of time in which the “independent states” have been going their separate ways since the collapse of the USSR, we find a direct correlation between movement towards communism and the development of society as a whole. This concerns not only, or even primarily, industry and the production of goods, but rather the opportunities for the development of the individual and of human abilities available to the broad masses of the population.

As early as 1925, there were 2,250,000 more pupils in primary schools than in 1913. The number of students in vocational schools doubled. Society went straight from the Civil War to the classroom. The right to an education was won for those who, before the Civil War, had been doomed to ignorance. Just think about it: the Soviet Union was only a year old, and it already had something to show the capitalist world, where the estate principle dominated education. In the capitalist world, a good education was, and still is, out of reach for the majority.

Within 20 years of October, the problem of child homelessness was solved, and the children who had lost their parents during the First World War and the Civil War were given not just a roof over their heads and bread, but an education, a trade, a start in life. Meanwhile, homeless children have reappeared amid the fragments of what was once a single socialist country—and most of them have living parents. The general level of education of the population has fallen and continues to fall. The education systems of the former Soviet republics are striving to conform to the new social realities: ever-narrower specialisation and a deepening divide between the educational opportunities available to the children of the poor and those of the rich. You will find a wealth of telling evidence of this process in Russia in the film “The Last Bell” by Konstantin Syomin and Yevgeny Spitsyn.

In the USSR, the campaign to eliminate illiteracy (likbez, launched in 1920) and the programme of universal primary education (vseobuch, from 1930) managed to solve the problem of primary education in less than twenty years—and this in an agrarian country where, before 1917, more than 80\% of the population was illiterate or semi-literate. By 1936, around 40 million illiterate people had been taught to read and write in the USSR. Between 1933 and 1937, more than 20 million illiterate and about 20 million semi-literate people studied in the officially registered likbez schools alone. According to the 1939 census, literacy among people aged 16 to 50 was approaching 90\%. It was these very people—not mythical warrior-heroes with crosses on their chests and God in their hearts (as Nikita Mikhalkov’s film “The Citadel” and other products of today’s propaganda industry would have it)—who defeated the war machine of Nazi Germany, the most powerful, modern and well-equipped of its day, which by the time of the invasion of the USSR had most of European industry working for it.

So, within 20 years, illiteracy was virtually eliminated and a network of medical facilities was built that gave everyone access to healthcare (doctors and teachers appeared in places that had never had them before); schools, libraries and cultural institutions sprang up, and science received a powerful boost, thanks in part to the creation of the necessary infrastructure. The building of communism, understood as prioritising the all-round development of the individual (education, healthcare, science, culture), made it possible to defeat fascism, which threatened the entire planet, and then, after a hard and devastating war, to rebuild the country in record time.

And what did life in the years after the breakup of the USSR—the complete dismantling of communist tendencies in social development—bring to the former Soviet republics?

In Ukraine alone, according to the State Statistics Committee, the absolute population decline from the collapse of the USSR up to 2014 amounted to almost 6.5 million people (before the collapse, population growth had been positive). This figure does not account for those who are still “registered” as citizens of Ukraine but have been forced to leave their homes to find work elsewhere. This in itself testifies to a colossal social catastrophe. Of course, someone will bring up the famine of 1932–1933. Mass famine is always a terrible tragedy, regardless of its true causes. But even according to the most anti-Soviet calculations (calculations, mind you, not groundless ideological claims) carried out by the Institute for Demography of the National Academy of Sciences of Ukraine, the famine in Ukraine claimed not 7 or 10 million lives, but 3.9 million. These calculations did not simply count losses; they compared the actual population with the projected population, taking into account the significant positive growth in the years before and after the famine. If the absolute population losses in Ukraine after the breakup of the USSR were calculated using the same method, the picture would be far more horrifying than the nearly 6.5 million obtained by simple subtraction.

Lithuania, Latvia and Estonia have also recorded significant depopulation. In the Russian Federation, thanks to migration, the decline was smaller: from 148.6 million in 1991 to 144.3 million in 2016.

If we take into account labour migration in other directions and the employment structure within the Central Asian republics themselves, it becomes clear that the overall quality of life, and the chance for citizens of these republics to find work at home that can support a family, have declined considerably.

Wars on the territory of the former USSR have become a frequent and commonplace occurrence: the conflict between Azerbaijan and Armenia in Nagorno-Karabakh (which began while the USSR still existed, but essentially during its collapse), the Civil War in Tajikistan, Transnistria, the two Chechen wars, Georgia and Abkhazia, South Ossetia, Donbass. And in all these wars people killed one another, in the last instance, to serve the interests of rival factions of capital—that is, their own exploiters.

Some may object that the Civil War that began in 1917 was also bloody and claimed, by various estimates, between 8 and 13 million lives on all sides. This figure includes not only those killed in combat but also those who died of disease, famine and terror. Yes, that is true, we reply. And, of course, it would have been better to avoid war altogether. But, first, it was not the Bolsheviks or the proletariat who unleashed it, but representatives of the old exploiting classes who did not want to lose their dominance. Second, in that war the working people were defending their own interests, not the interests of those who live off the exploitation of their labour. Third, that war was a struggle for the opportunity to carry out a radical transformation of the whole of society, not merely a redivision of capital’s spheres of influence. The Great Patriotic War deserves to be discussed separately, since it was a struggle against fascism in which the Soviet people defended not only their own interests but those of all humanity.

As you can see, even such a simple comparison of the first decades of the USSR with the life of the republics after its collapse allows us to speak of the advantages of a society in which there is room for the tendency towards communism. And this despite the fact that the process was far from smooth, and there was much that could by no means be called communist. We are referring to purely capitalist processes and phenomena within socialist society. It is precisely these that critics of the USSR make so much noise about, even though this criticism, where it is entirely justified, turns out to be self-criticism with respect to the capitalist world.

The tendency towards communism in the USSR made it possible, just twelve years after the Great Patriotic War, which claimed millions of lives and destroyed a substantial part of the country, to launch the first artificial Earth satellite—and soon afterwards to send the first human into space. The building of communism made it possible to work out scientifically ideas that, for us in the present state of society, are ideas of the future rather than of the past, since their realisation is incompatible with private property in the means of production. The list could go on at length, and it could include not only Soviet achievements but also those of other socialist countries—countries that arose precisely along the path Lenin envisioned in the light of the hypothesis of a victorious world revolution.

Of course, this grand construction project involved not only achievements but also retreats and failures. But even if we set aside what has been said above, the argument

“we failed to build communism, so it is impossible altogether and nothing will ever come of it” suffers from the same logical flaw as the claims made by opponents of the ideas that underpinned, for example, the emergence and development of aviation. Those opponents insisted that it was impossible to build flying machines heavier than air. It is the same as saying that a person, or humanity as a whole, will never be able to learn something simply because they have not managed it so far, even though attempts have been made.

Such illogical arguments involve no examination of the essence of the matter and no substantive objections. That is why they have nothing whatsoever to do with the science of communism. The failure of humanity’s first leap towards communism testifies only to the need for a detailed study of historical experience in order to understand the reason for that failure, and for a more thorough study of the conditions that must be achieved for the transition to communism. The very existence and content of the communist ideal, which was born as an understanding of the tendencies and contradictions of capitalist society, demands this.

So, since communism, among other things, has become a science, let us proceed to study it. This pamphlet, of course, cannot claim to be a detailed and comprehensive study of communism. Rather, it is a text intended to set out the basic tenets of the science of communism as developed in Marxist theory. A deeper and more detailed treatment of the problem is a task for the near future. We invite our readers to join in this work, including in the form of criticism of our own.
\ifSubfilesClassLoaded{\setcounter{page}{14}}{}
\section{The Foundation of Communism}
\label{sec:03_foundation}

The foundation of capitalism is commodity production. All the wealth of capitalist society appears as an aggregate of commodities. Since the products of labour under capitalism inevitably take on the commodity form, the commodity is the “cell of capitalism”.

In accordance with the materialist conception of history, the foundation of communism as a social order can likewise only be a corresponding mode of production.

All the wealth of communist society lies in the universality of the content and forms of human activity. Marxism in general sees the essence of man in activity aimed at transforming nature (including his own social nature). To labour, to work up natural material, to create things that do not exist in nature—this is what fundamentally distinguishes human life from the life activity of an animal and determines all the other distinctions. Yet until now labour itself, on a social scale, has been unfree. Every pre-communist mode of production was marked by its own specific form of the unfreedom of labour. Communism, by contrast, presupposes the complete overcoming of this unfreedom. Therefore its foundation can only be free labour “irrespective of quotas and without reward”, as man’s “life’s prime want”.

The simplest directly collective form in which such labour is carried out (or, better said, such activity, since the word “labour” still retains a negative connotation to some extent) is the very “cell” of communism, which contains within it the further development of communist society in all its necessary determinations.

Free activity and its directly collective form are a reality in our time, just as commodity production and the commodity were a reality even before the emergence of capitalism. And the transition to a new social system will take place when man’s free activity becomes a universal relation, just as the commodity once did. From the motley array of commodities on the market, each with its own properties, one cannot directly grasp what a commodity and commodity production are. Likewise, from the diversity of types and methods of organising communist activity in their particularity, one cannot grasp what it is and what its directly collective form is. All the more so because words can obscure the essence of the matter. The directly collective form of free activity can exist even when a person carries it out entirely alone. We must single out those determinations that are necessary to it.

In the directly collective form, free activity is carried out by the measure of freedom, truth, goodness and beauty, and hence of the highest human interest; that is, it is a form of sensuous activity valuable in and of itself. Free activity necessarily presupposes the unity of all three components—sensuousness, morality and thinking—and their elevation to the level of freedom. In the course of such activity, man is affirmed in all his fullness, regardless of what exactly is being done and of which element (truth, goodness or beauty) comes to the fore.

This directly collective form is merely the activity itself, together with its corresponding content, viewed from the side of form. In everyday speech, such activity is often called social creativity, or simply creativity. Form matters because development takes place and finds expression in the necessary movement of form. Content always calls into being a form corresponding to it and is realised through that form. That is why the contradictions of content are first grasped by consciousness as contradictions of form. It is in the unfolding of form, and through that unfolding (the succession of forms), that the movement of content actually takes place. On the side of form, the important aspects of activity are the nature of the ties between people, the goal, and the social being of the object and means of labour.

Direct collectivity in free activity is expressed, first and foremost, in its orientation towards other people. In directly collective forms of activity, the human being always appears as the goal and the highest value. Everything that is done is done directly for the development of other people. For this reason, the product of such labour never acquires an independent thing-like form (the form of a commodity). The product of such labour cannot confront living people as a thing alien to them. In other words, this activity is directed towards the development of society and the improvement of other people’s lives; it is carried out for them, and therefore also for oneself as a member of that society. What are now customarily called material values can serve only as means in free communist activity.

As for the nature of the ties between people, it can, by and large, be defined as comradeship, or, what amounts to the same thing, free association.

The nature of these ties also points to the means of such activity, since in a free association none of the participants can serve as a means for carrying out someone else’s will or interests alien to him. Even if someone, thanks to his abilities and skills, takes the lead in the collective, the others do not blindly carry out another’s will; rather, they cooperate in accomplishing a common cause that is equally important and necessary for all its participants.

The nature of these ties also points to the social being of the means and subject of communist activity. Thus, only things and processes distinct from human beings can serve as means. The social form of being of the means and subjects of labour presupposes their directly social use. It was in precisely this sense that Marx spoke of the transition from the government of persons to the administration of things.

From what has been said above, it is clear that such forms of human activity are free of any external coercion whatsoever, whether direct or indirect (remuneration for labour), since neither the process nor the result of communist activity is a matter of indifference to the one who carries it out. In each instance, a person has a direct interest in his own labour, unlike the situation that is normal for modern society, in which the activity itself serves merely as a means of obtaining money.

Of course, direct collectivity also consists in the fact that the labour of any member of society rests on the labour of other people, that the subject, means and method of labour are products created by other people, while his own labour, in turn, is an integral part of all social labour and provides material for future activity. But this is characteristic of any labour carried out in any social form, not only of communist activity, and therefore it does not concern what is specific to communism.

By the directly collective form of free activity we mean the association of people in the process of activity, where the goal of each is the development of all, and the means of achieving it are things and natural processes with respect to which free and equal expression of will takes place.

The directly collective form of activity is not merely an external social form of activity, one that may or may not be present without affecting the content of the activity; on the contrary, it is its substantive form. Outside it, free activity by the standards of truth, goodness and beauty cannot be carried out at all.\footnote{What is meant here is by no means the common representation in which beauty is associated with taste, and goodness with opinions about what is good and bad, which may vary, while only truth is accorded any kind of objective significance. The categories of “truth”, “goodness” and “beauty” were developed over millennia as objective ideal characteristics of human activity in the world, unfolding (developing) within that activity, and therefore also as objective characteristics of the world (of matter). In other words, to put it briefly, when we speak of activity by the standards of truth, goodness and beauty, we are talking about the “points of growth” of human sensuousness, cognition and morality, about the development of the relationship of human to human and of human to nature, taken from the side of the ideal dimension.} The movement and metamorphoses of such a form presuppose the preservation of the aforementioned moments, just as value is always preserved in the metamorphoses of the commodity form of the products of labour (that which makes a commodity precisely a commodity, rather than simply a product of labour).

As for the content of such activity, it is always revolutionary. Herein lies its contradiction with respect to itself. Free communist activity constantly posits both its own internal limit and the transcendence of that limit. By means of historically attained, and therefore historically limited, frontiers of activity, this activity itself changes, positing both new content (thanks to the widening horizon of the assimilation of matter) and a new form. This does not happen in a roundabout way, independently of the participants in the labour process, as it does in the production of things. On the contrary, this change happens purposefully. Therefore communism, as man’s assimilation of his own activity, lays bare the internal contradictoriness of human activity in general, and hence the internal contradictoriness of the material world in which man carries out his activity. That is why Marx’s remark in the first of the Theses on Feuerbach is of paramount methodological importance for the science of communism. “The chief defect of all hitherto existing materialism—that of Feuerbach included—is that the thing, reality, sensuousness, is conceived only in the form of the object or of contemplation, but not as sensuous human activity, practice, not subjectively. Hence it happened that the active side, in contradistinction to materialism, was developed by idealism—but only abstractly, since, of course, idealism does not know real, sensuous activity as such. Feuerbach wants sensuous objects, really distinct from the thought what is good and bad, which can vary, while only truth is accorded some kind of objective significance. The categories of ‘truth’, ‘goodness’ and ‘beauty’ were worked out over millennia as objective ideal characteristics of human activity in the world, unfolding (developing) within that activity, and therefore as objective characteristics of the world (of matter). In other words, to put it briefly, when we speak of activity by the standards of truth, goodness and beauty, we are speaking of ‘points of development’ of human sensuousness, cognition and morality, of the development of man’s relation to man and of man’s relation to nature, taken from the side of their ideal dimension. objects, but he does not conceive human activity itself as objective activity.” This remark does more than emphasise the role of idealism in the becoming of the theory of activity—although that role is indeed enormous and important, so much so that Lenin later spoke of the importance of developing societies of “materialist friends of Hegelian dialectics”. The categories of dialectics, understood materialistically as “forms of human activity repeated millions of times” (Lenin), together with their internal contradictoriness and movement, are an indispensable tool for developing the theory of communism precisely because they reflect the essential moments of human activity in general and capture that without which human activity cannot, in principle, take place. But Marx’s thesis also places emphasis on the materiality of human activity.

Let us emphasise once again that free activity is not “leisure” in the modern sense of the word, not idle pastime, but activity by the measure of freedom. Philosophy, in the person of Baruch Spinoza, once defined freedom as recognised necessity. True freedom is defined by acting in accordance with this necessity.

In their free time, in accordance with historical necessity, people develop culture in the broadest sense of the word. Such activity is not something invented by communists in fantastical dreams. It existed in the past and exists now. But it existed, and continues to exist, within very rigid limits and is the privilege of only a few. The overwhelming majority of the members of society do not know free activity. This has been and still is the case, among other reasons, because one needs not only to have the time but also to be developed enough to be capable of carrying out such activity. The conditions for such development have been, and remain, highly limited in the modern society of commodity production.

So, in order to move in the direction of communism, what must be produced is communism itself, that is, free activity, or more precisely, free time and a corresponding level of development of the members of society who could carry it out.

Consequently, distribution here is, first and foremost, the distribution of free time and of opportunities for the assimilation of human culture. The distribution of the products of labour is necessary only insofar as it provides free time. It is in this spirit that the question of the efficiency of social production in the period of transition to communism should be posed. Under capitalism, free time is a cost of production; it is not accounted for, not valued, and therefore often never becomes truly free and is not used for its proper purpose. The purpose of free time is to be used for the continuous development of social relations (or the development of activity, which is essentially the same thing).

That is why, for communism, it is very important that the habit of treating one's own free time, and the activity carried out during it, as social wealth not subject to private squandering should become the common possession of all. In general, the force of habit and the collective formation of habits are of enormous significance here, if only because being free, that is, acting by the measure of freedom, is always very difficult. It means being at the cutting edge of the contradictions through which the development of society takes place. This is not only hard work but also a tragedy, in the sense that one must constantly overcome oneself and negate one's past self. And collective social freedom is a collective social tragedy and its constant resolution. Communist society is not a society of universal laziness and idleness, but a society of universal free activity, which is for people life’s prime want.

Communism is realised through the necessary succession of directly collective forms of human activity aimed at transforming that very activity (at developing people’s life itself). The forms already attained, striving beyond their own limits, give rise to new ones. That is why this succession cannot be pictured as the same metamorphoses endlessly repeating one after another (as we see in the movement of capital). The inner logic of the movement of the directly collective forms of free activity, the necessity of their succession, is determined not merely by the historically attained mode of production and the relations of production corresponding to it, but by freedom as the necessity of matter itself, which unfolds as matter is assimilated.

In its most general form, the law of communism can be formulated as the law of freedom, or the law of the self-determination of human activity: the free activity of the members of society, carried out as the conscious development of the content and forms of that activity, determines the development of society as a whole. Or, in other words, freedom, in realising itself, strives towards self-expansion.

Of course, understanding how this law operates requires constant refinement, drawing on all available knowledge of both nature and society. That is why, under the appropriate conditions, the theory of communism can play a unifying role for all the sciences, pointing to the real connection between them as the connection of human practice itself.

The directly collective forms of free human activity that exist today as “shoots of communism” are not merely evidence of its possibility; they are its real being. It is important to understand that the communist mode of production already exists; the point is only that it should become dominant and extend to the whole of society.

In their movement, the directly collective forms of free human activity, as “shoots of communism”, can unfold into communism as a social structure, that is, as society’s mode of existence. In this sense, one of the most valuable propositions in Lenin’s theoretical legacy is his advice to seek out and nurture the “shoots” of communism wherever they are found.
\ifSubfilesClassLoaded{\setcounter{page}{19}}{}
\section{Communism as a Reality in Social Production}
\label{sec:04_communism_reality}

The claim that communism already exists may seem too bold. At the very least, it requires substantiation. What strikes the eye is the opposite: the absence of communism, the dominance of labour whose value to the one who performs it is determined solely in the form of payment for labour, that is, wages. It may seem that the cell (the germ) of communism, the directly collective form of free human activity, is something that does not lie on the main road of the development of human society. We will try to show, however, that this is not so. To do this, it is not enough to select examples that would attest to the presence of communism within the truncated confines of exploitative class society. We must examine the directly collective form of free activity produced in contemporary class society from the standpoint of the possibilities for its further unfolding.

Recall that in his time Lenin was able to see first the Soviets, and later the communist subbotniks, in this light. “The communist subbotniks organised by the workers on their own initiative are really of enormous significance in this respect. Evidently, this is only a beginning, but it is a beginning of exceptionally great importance. It is the beginning of a revolution that is more difficult, more tangible, more radical and more decisive than the overthrow of the bourgeoisie, for it is a victory over our own conservatism, indiscipline, petty-bourgeois egoism, a victory over the habits left as a heritage to the worker and peasant by accursed capitalism. When this victory is consolidated, then and only then will the new social discipline, socialist discipline, be created; then and only then will a reversion to capitalism become impossible, will communism become really invincible.” Lenin maintained that the first communist subbotnik, organised on 10 May 1919 by railway workers of the Moscow–Kazan Railway, was of greater historical significance than any victory in the First World War. The point was not merely that the workers who organised the subbotnik had expressed their own will, but also that this undertaking brought a substantial increase in the productivity of labour, owing to the fact that for each worker it was a personal matter, one in which their personal interest in developing the railway economy for the good of society as a whole was realised.

After we published the chapter on the foundation of communism, one of our readers made a crucial remark: “The communist mode of production also implies the production of material goods, and not simply some Pushkin literary circle that fosters the becoming of poets. Because ‘communist relations’ based on exploitation for the ‘chosen few’ are plentiful today,” the reader asserted. He demanded that we present evidence that communism already exists, and show what it could be like specifically in the sphere of material production. In our view, the reader is right on two counts. First, the reader spoke on behalf of the proletariat—people directly engaged in material production through the sale of their own labour power, where no communism is to be seen. Second, the reader raises the crucial question of communist material production, since it is precisely this that determines communism on a social scale, and not the aspirations, wishes or even actions of political revolutionaries.

However, communism is impossible in purely material production. Where a person directly produces only things, while the production of social relations takes place in this process entirely independently of him—of his aims, will and consciousness—there can be no talk of communism whatsoever, even if the production of things in itself gives the worker a certain pleasure in labour. In such production, a person is merely a material force alongside other material forces. In the last instance, he acts as a thing affecting other things. From the standpoint of Marxism, for communism to arrive, such labour not only can but must be abolished as human labour and handed over entirely to machines. Speaking of the historical role of capital, Marx specifically noted the importance of creating surplus labour and the universality of human needs as conditions and prerequisites of communism. Capital already manages this even under capitalism. But this alone is not enough. The historical role of capital will be finally fulfilled, and communist society will become a reality, only when “labour in which a human being does what he can make things do for him, for man, has ceased”. In this sense, a person cannot act in a communist way until he has been taken beyond the bounds of purely material production, so long as he does what things can do and, importantly, himself acts as a thing in this process. That is why the proletariat is the only class that is consistently interested, to the very end, in seeing the historical role of capital fulfilled completely. This is precisely the abolition of the conditions in which society is constantly split into classes.

The author of the comment rightly pointed out that in the sphere of the production of the ideal one can find as much communism as one likes. This includes the work of a scientist and of a collective of scientists who labour day and night to give people scientific truth, and on top of that give educational lectures in their free time. A certain minimum of material goods is here a condition, not the goal, of engaging in communist activity. Communist activity is the work of the writer, the artist, the musician who labour so that people may see, hear, feel and understand the world better, so as to push back the boundaries of human abilities, and not for the sake of money or fame, and so on and so forth. Of course, by no means all activity in the sphere of the ideal is communist. But if the features of the directly collective form of free human activity outlined in the previous chapter are present, then we are dealing with local communism, even though communism exists here in a very truncated form.

Communism here, as we will rightly be told, is far from complete. First, the very conditions for it are produced at the expense of purely capitalist exploitation of man by man. Second, capital, as the dominant social relation, confines freedom within rigid limits here as well. It is impossible to build communism as a mode of material production without going beyond the bounds of the sphere of the ideal. And although communism in the sphere of the ideal has played a crucial role throughout human history, since it was in free activity that the ideas were born which do not merely reflect but also “create” the world, it still cannot be of decisive significance, since it does not yet make society communist. Third, there is as yet no question here of overcoming the social division of labour, which is the most important condition for communism. The universality of man is possible here only at the level of consciousness, not at the level of the wholeness of man as an objective-practical being.

However, there exists a kind of material production that is not purely material production but is the production of man in his wholeness. Such production does not merely determine the production of the ideal in a roundabout way, but incorporates it: the production of social consciousness, of the social ideal, of social goal-setting. The example given above of the communist subbotnik, during which workers loaded and unloaded freight and repaired rolling stock, illustrates this well. Unappealing at first glance, the manual labour here was not merely the labour of carrying heavy loads but labour aimed at producing a new social reality. It is fitting to recall here the subbotniks for harvesting sugar cane in Cuba: with fairly primitive tools, under a scorching sun, the participants were not simply cutting cane but earning the means to modernise production and build schools and hospitals, in order to create the conditions for the development of every member of the Cuban people. Such labour also existed in the Gorky Colony and the Dzerzhinsky Commune, both led by Anton Semyonovich Makarenko. Any task, even digging cesspits with a shovel and building wooden outhouses over them, became here, in the last instance, the production of the new man, free human activity by the standards of truth, goodness and beauty.

The question of the free form of directly collective activity in material production must therefore be posed as follows: what connection with society, as a whole in its development, does such activity embody and produce for the person who works, and what place does it occupy and what role does it play in the development of the social whole?

It is easy to see that the examples given above are taken from the history of socialism. Under socialism, the very fact of the formal socialisation of the means of production removes a number of obstacles to the free transformation of social life. We may therefore be told that these examples are not entirely valid. But, first, with the sublation of barriers, what gains the opportunity to be realised is not something that does not exist at all, nor something for which there are no conditions whatsoever, but something for which the full set of conditions was lacking—and the only thing missing from that set was the transcendence of the rule of capital. The fact that the sublation of barriers gives free rein to the development of communist tendencies means, at the very least, that these tendencies exist, and that the potential for the becoming of such activity already lies hidden within capitalism.

It may also be objected that the subbotnik is something unstable and situational, something that does not determine the course of life as a whole. Although in the course of a subbotnik social goals are set and achieved as the personal goals of its participants, outside the subbotnik the same old capitalist laws continue to operate in material production: labour power is a commodity, and each individual proletarian enters into a contract for the purchase and sale of labour power with his own class, which confronts him as a collective capitalist; the law of value retains its force unless its operation is deliberately restricted by pushing capital’s tendency towards the automation of production to its limit.

In other words, what has been cited above as an example does not, by itself, make society communist. These remarks are perfectly correct, and we shall not dispute them. But the essence of the matter lies not only, and not so much, in simple registration, in merely noting that this is communism while that is not communism, or not quite communism, but in those tendencies of society whose potential is contained in, and can be seen in, such islands of communism. It was precisely in this spirit that Lenin thought about the “shoots of communism”. He was interested in subbotniks precisely as shoots, as germs. That is why he draws attention chiefly to two things: 1) the free expression of will by the participants in the subbotnik in relation both to the means of production and to themselves; 2) the significant growth in the productivity of labour during subbotniks, the achievement of which is bound up with this free expression of will.

Another significant remark would be to point out that in the examples given above, the activity is carried out with tools that are far from the most advanced, even by capitalist standards, and certainly not on a communist technological base. This straining of human effort stems, rather, from a lack of social development: a lack of technology, a lack of organisation of technological processes, even in purely capitalist terms. In our view, this is the most reasonable objection to such communism as the germ of a new society. The communism out of need discussed above is not yet communism out of inner necessity, and therefore the development of a communist society does not yet follow from it on its own. It is only a precondition for such a society, albeit an essential one. Here, once again, we need to investigate the connection, the continuity in terms of the ways collectivity is built, between this kind of communism and communism developed on its own foundation, a task to which a separate study must be devoted.

However, pointing to the need for a separate study of this problem by no means relieves us of the obligation to show, in general terms, the practical possibility of this connection, nor of the obligation to provide evidence of the existence of directly collective forms of free human activity based on advanced modern technologies.

The connection between “communism out of need” and communism on its own basis, their continuity, may consist, first, in the upbringing of the masses, in the formation of a habit of free expression of will in accordance with social necessity in the production of social life; and second, in the creation of communist forms of collective organisation which, once the social need that called them into being has been overcome, need not necessarily disappear, but, given the upbringing they foster in the spirit of genuine freedom as conscious social necessity, may become forms of the communist organisation of social life. We are not utopians. We shall not invent exactly how this will come about. That depends each time on the interweaving of concrete-historical circumstances. We merely note the possibility of such further development, inherent in such shoots as regards the collective-organisational form of human activity, which presuppose free and equal participation in it.

As for communism based on advanced technologies, we see such shoots in the joint use and creation of modern information technologies. For example, the Linux operating system and other programs developed within the framework of free software production have been created on communist principles.

We should make a caveat: we are not talking about all programs with open source code. Many large modern projects have given rise to huge corporations, and the mere fact that the code is open does not make the activity of creating programs free. Where this is the case, the corporation completely determines the project’s policy, sponsors it and profits from it. The code is written by programmers for pay. There is no communism here (although there is a possibility of transforming such activity into communist activity). This is how, for example, auxiliary products needed by several companies at once are written. Cooperation makes development cheaper. But in the last instance, the product here is produced for profit. Often, any use of such a product at all, let alone commercial use, requires large resources that only big corporations possess. They usually have a commercial service or software for which what has been developed is one of its “components”, so to speak, that is, a condition for its operation and successful sales. But the corporations have monopolised something that technologically has serious potential, going beyond the bounds of capitalism. This potential is tied to the technical possibility of both producing and using software, thanks to the development of information machines and networks.

Free software began, and some projects now exist as projects in a space where there are no commodity-money relations. Moreover, developers are waging a struggle to ensure that no one derives commercial benefit from this. This is precisely empirically given communism within the framework of capitalism. Open-source software presupposes full free access to programs and to their code, as well as making freely available the programs themselves and their code created using previously obtained free source code. For example, the Linux operating system proved competitive with commercial Windows. Corporations ultimately try in every way to bring it under their control, and they build products based on the open source code of this system. But the system itself still remains free of charge.

Admittedly, what is created within the framework of non-commercial free software by no means always works as intended, as users point out. But commercial software also goes through this on its way to becoming a fully finished product, and some developments are abandoned altogether because they lead to dead ends. The same is true here. The only difference is that here all stages are open and visible to everyone. What is unfinished and what has hit a dead end also become public knowledge. Free programs are in general often a semi-finished product, designed to be brought to full readiness by the user. This hampers their competitiveness under capitalism in a world with a generally low level of information literacy.

Corporations develop their products with this level of literacy in mind, refining them to the point where they can be used in a primitive, intuitive way. Their product is intended for people who, under conditions of alienation, find it easier to pay for intuitive usability than to spend time figuring out a program. The corporations' product does not presuppose that users will create new programs or modify existing ones to suit their needs by adding to or rewriting part of the code. Using them requires no particularly high level of culture, let alone an understanding of technological processes. That is why programs produced by large corporations have far broader sales opportunities. Moreover, they are often of higher quality, since they have been worked out in every detail.

Free software, non-commercial in every sense, by contrast, demands at least a basic grasp of programming fundamentals; it presupposes the possibility, and often the necessity, of extending and rewriting it. It compels users to become programmers themselves, at least in part, and to keep improving their skills\footnote{Granted, this can fray your nerves and make it hard to focus on the task at hand. In the end, instead of getting the job done, people have to keep fixing a Linux tool that is constantly falling apart. But there are no fundamental obstacles (from the standpoint of potential) to overcoming this and writing better-quality programs, provided that more social time can be devoted to creating them.}. It is precisely this last circumstance, which limits the purely capitalist competitiveness of free software, that points to its greater social potential compared with commercial software (this is not about giving up a convenient interface, but about being able to create a convenient interface oneself). Built into its production, its means of distribution and its use is not merely a potential for technical development, but a potential for human development: for overcoming the division of labour and bringing all participants in the process up to the level of contemporary computer literacy. And although capitalist companies try (sometimes successfully) to bring this local communism under their control, it can quite reasonably be counted among the shoots of a new society, however weak and tender they may still be. First, it arose on the basis of the most advanced technologies capitalism has to offer. Second, it draws a clear line between the capitalist way of using and creating these technologies and a way that goes beyond the bounds of capitalism.

It is also important that such production, on purely technological grounds, already leads to the abolition of hierarchy in production and of the separation between users and programmers. It requires no external link between them, nor any special social “governing bodies” to provide that link. The external government of persons, even in its purest democratic form—that of deciding questions by general vote—is no longer needed here. It tends to wither away, becoming superfluous.

It may be objected that, in speaking of the non-commercial, free production of software, we are once again dealing with the sphere of the production of the ideal. Indeed, the process of writing program code, as a kind of text, belongs to ideal activity, and program code (though not the program itself as it ordinarily runs) is an ideal product. However, this is the kind of ideal activity that forms a vital part of modern industrial material production. Creating software is the labour of introducing automation into industrial production. And although program code is an ideal product, the operation of modern industrial machines based on these programs (that is, the “work” of programs in industry) is purely material machine production.

Of course, the production of free software also yields a great many pointless and useless programs. That does not really matter. What matters is that it produces, and can produce, not only these, but also products fully fit for use on the scale of modern industry. And these products are produced communistically. Often, though, such projects are either never finished or drag on for years, because the participants have to earn a living and therefore have relatively little time for free activity.

We see the same thing in the creation and use of modern databases. Again, far from all of them are produced communistically. But there are large publicly accessible databases that presuppose exclusively joint creation and use. They simply cannot be produced any other way. Otherwise they would become pointless and unfit for the very purposes for which they came into being. For example, some spatial databases presuppose free participation in their formation in order to build a picture of space in the interests of society's development.

Here, too, the objective laws of human activity tend to exclude the mandatory hierarchy characteristic of the capitalist mode of production, drawing all participants into the process as equals and as free individuals.

The structure of a database is also an ideal product, we will be told. But in the creation of databases, ideal activity is so fused with material production (including material reflection by means of databases) and with the reproduction of social life that, for all its ideality, it certainly cannot be classed as superstructural activity. As in the case of software, it is an ideal moment of material social production, one that has not been abstracted (not in consciousness, but in reality) from that production in the form of separated-off forms of social consciousness.

Thus, what we are pointing to here is not merely ideal activity, but industrial material production, and the tendencies and potentials already contained within it under capitalism, which, once the barriers in the form of private property in the means of production are sublated, can develop into communism, forming both a purely technological and a collective-organisational basis for it, one that educates and creates the new man.

The further development of this tendency leads to the automatic factory in its idea, in which Marx saw the prototype of future material production. Here, too, lies the potential for realising V. M. Glushkov’s ideas of “nationwide” (society-wide) systems of economic management

(of this automatic factory on a social scale). In the long run (provided that social barriers are sublated), everyone can have access to all information on the state of all social production in real time. Everyone can take part in managing not just the production of things (the production of things makes sense only insofar as it serves the main goal of production), but the production of man—of social life in its entirety.

What is interesting is that participation in all these processes presupposes the production and consumption by all participants of an ideal product, but one that is not alienated and cannot be alienated from material production. In this we can trace the emergence of a new type of social consciousness, one that sublates its alienated forms.

And the need to sublate the social barriers to free activity is likewise created by modern capitalism. Independently of anyone’s will or consciousness, capitalism itself objectively puts this question on the agenda. The objective logic of the movement of modern capital constantly gives rise to wars and thereby makes people’s customary modes and forms of activity impossible. Capital forces people to create qualitatively new forms of activity in order to solve the problems of the reproduction of social life while resisting the destructive power of capital.
\ifSubfilesClassLoaded{\setcounter{page}{27}}{}
\section{Common (Social) Property}
\label{sec:05_common_property}

The obstacle to the movement towards communism is private property. The first work in which Marx gives serious and thorough consideration to communism is the “Economic and Philosophic Manuscripts of 1844”. The chapter on communism appears in the second manuscript, immediately following the chapter on private property. It opens with reflections on private property, on labour as the essence of private property, and on how private property and labour are understood by various representatives of socialism and communism as an ideological current. Communists (including those before Marx) oppose private property in the means of production because it makes people the slaves of other people and gives rise to all the horrors of capitalist society.

Private property and personal property are not the same thing. Items of direct use and consumption belonging to their owner are, in form, private property, just as plants, factories, shops, corporations and large tracts of land are. This circumstance obscures the essence of the matter. That is why, when anti-communists hear about the abolition of private property, the most primitive among them think it means the socialisation of absolutely everything and start fearing for their personal belongings. On this point, we can reassure them. Communists are interested in articles of personal consumption not in the sense of depriving anyone of them, but only in ensuring that everyone is provided with them to a degree sufficient for a full and many-sidedly rich life. Economically, these items are not private property but personal property.

The abolition of private property does not do away with personal property (private property is a relationship between people concerning things, not the things themselves), since consumption must necessarily be personal in character. On the contrary, the abolition of private property ensures that all members of society have access to the items they need for personal consumption, and, where it is more rational to do so, it shifts consumption into social forms. For example, organising consumer cooperatives makes it possible to use things more efficiently and reduces the need to produce and store them, and therefore also the amount of social time spent on this. Thus, organising a simple consumer cooperative for sharing items that are not used very often makes it possible to avoid cluttering up space and to use each item as efficiently as possible right up to its moral and physical depreciation. The item passes from hand to hand, and every member of the cooperative can derive a useful effect from it. Examples of consumer cooperatives have been described in the literature many times. This is easily done with a wide variety of items, which cease to belong to anyone personally while still being used by everyone to the extent needed to satisfy their needs. Here, communism stands opposed to personal property not at all from the side its opponents fear. Social forms of consumption are aimed not at taking something away from someone, but at giving everyone access to articles of consumption.

As for such social forms of everyday life as the shared use of washing machines, communal cooking and food delivery, and the shared use of cleaning equipment, the savings in resources and human time (both directly and indirectly—as savings in the time spent producing those resources) can be very large while maintaining the same level of consumption. At any given moment, these savings are proportional to the capabilities of the technology. But even now, for example, cooking porridge for one child at home takes the same amount of time as cooking porridge for 200 children in a kindergarten or a milk kitchen equipped with modern appliances.

The same applies to cooking for adults. Social consumption makes it possible to use more efficient technologies that cannot be applied in individual households, with a much higher level of automation in every process. At the same time, quality control can be carried out better, which ultimately benefits everyone involved. Movement in this direction exists under capitalism as well. Capitalism itself, in its own way, produces and standardises systems of public catering (McDonald's, for example, offers identical French fries in every part of the world), laundries, and so on. Under capitalism, however, they are still not used in a social way. Forms of collective consumption constantly run up against resistance in commodity production, and production, as we know, determines consumption.

All the burdens of everyday life, in which a great many people are engaged, can be replaced by advanced production based on modern technology; private transport, which gets in its own way and that of pedestrians and creates traffic jams, is replaced by a well-developed public transport system, and so on, in every sphere of human life where this would benefit everyone. For example, in today's large cities with developed public transport, those who use it, especially at rush hour, reach their destinations more efficiently. In Moscow, children who take the metro to school mostly arrive on time, whereas those who are driven in private cars tend to be late, because attempts to use limited public urban space in a private way prove less efficient than its social use. This also explains why city authorities, during snowfalls, urge people to give up their private cars and switch to public transport in order to minimise accidents on the roads.

Commodity production cultivates individual forms of consumption. For all their irrationality and inefficiency, they are a condition for constant demand for commodities on the market, a condition for the movement of capital. It is precisely in this sense that private property comes into contact with personal property and influences it.

Private property, or, to put it more precisely, private property in the means of production, is above all the ability to extract surplus product (surplus value, in the form of profit or in the form of rent), that is, the exploitation of other people.

Communists speak of the socialisation of private property, not personal property—of the means of social production, which today are used not for the benefit of society as a whole but serve the exploitation of man by man. Private property is not things, but social relations mediated by things. Its essence is exploitation.

Representatives of communism before Marx understood the abolition of private property as its extension to the whole of society, the transformation of all members of society into workers, and the equalisation of everyone with respect to both labour and consumption. Many representatives of pre-Marxian communism regarded private property merely as material wealth that could be owned by everyone in common. Proudhon, for example, conceived of private property only as capital, which was to be abolished. Closer to Marx were those theorists who saw the source of private property and of its pernicious influence on man in a particular character of labour, “labour which is levelled down, fragmented, and therefore unfree”, yet they reduced such labour to particular kinds of labour. “Fourier, who, like the Physiocrats, again regards agricultural labour as at least the best kind of labour, whereas for Saint-Simon, on the contrary, the essence of the matter lies in industrial labour as such, and he accordingly aspires to the exclusive rule of the industrialists and the improvement of the workers’ condition. Finally, communism is the positive expression of the transcendence of private property; at first it appears as universal private property.” “Positive” here means “preserving all the wealth of culture and the active capacities of man developed in a society of commodity production”.

Marx subjects this understanding of the negation of private property as its extension to the whole of society to merciless criticism. He shows that such communism does not overcome private property but is merely the theoretical culmination of private property as a social relation. He notes that the dominance of property, as an aggregate of things, over such communism, both as an ideological and theoretical position and as a social movement, is still very great. Social relations are still conceived only in terms of private property. Such communism “wants to destroy everything which is not capable of being possessed by all as private property. It wants to disregard talent, etc., in an arbitrary manner. For it the sole purpose of life and existence is direct, physical possession. The category of the worker is not done away with, but extended to all men. The relationship of private property persists as the relationship of the community to the world of things. Finally, this movement of opposing universal private property to private property finds expression in the brutish form of opposing to marriage (certainly a form of exclusive private property) the community of women, in which a woman becomes a piece of communal and common property.”

For Marx, the idea of the community of women is a relation to a human being as a thing, which also reveals man’s relation to the world of things created by man for man: “Just as woman passes from marriage to general prostitution, so the entire world of wealth (that is, of man’s objective substance) passes from the relationship of exclusive marriage with the private owner to a state of universal prostitution with the community.” Marx calls such communism crude. Wealth here means the entire aggregate of things created by people for people, in which are embodied both the active capacities of human society and the needs it has developed.

The theory of crude communism, which negates the human personality, merely absolutises in thought a real tendency of the society of private property. Where the wealth of human relations is reduced to relations of private property, the personality is reduced to an economic agent of private property. Indeed, the personality of these agents (the proletarian and the capitalist alike) is a matter of complete indifference to the movement of private property. The capitalist is interested in the proletarian not as a personality but solely as the seller of labour power of a certain quality, and the proletarian is interested in the capitalist solely as its buyer. The same thing happens in every act of buying and selling any other commodity. All the other personal qualities of economic agents have no bearing on the movement of private property. If people spend their entire lives without ever stepping outside these relations, their personality begins to be reduced to these determinations alone.

Crude communism presupposes a transcendence of private property that would not be a genuine assimilation of it. This is evident from its abstract negation of the achievements of culture, its treatment of them as superfluous merely because the proletarian, as a proletarian, has no need of them, and (in Marx’s words) its striving for the unnatural simplicity of the poor and crude man who has few needs, who cannot overcome private property because he has not yet even reached it.

The slogan “Take everything and divide it up”, typical of present-day exponents of pseudo-communism, stands lower still. They have not even reached socialisation in the form proposed by the exponents of crude communism, namely the socialisation of poverty.

Against the crude communist theories, Marx sets communism as the positive transcendence of private property. Such communism, as a historical process, preserves the wealth of all preceding historical development in the form of the active capacities, needs, and material and spiritual culture of humanity. Communism makes all of this the possession of every member of society. As a historical process, it is conscious rather than spontaneous in character, since socialisation presupposes a conscious attitude on the part of the members of society towards the production and reproduction of their social life. Communism presupposes the assimilation of the objective laws of social development, which humanity can put to its service in exactly the same way as it does the laws of physics, chemistry and biology in modern industry.

The whole preceding history of humanity creates the conditions for a communism that overcomes private property while preserving, in this very overcoming, all the wealth of preceding development. Over the course of all preceding history, a corresponding consciousness also takes shape: the theory of communism. Communism is understood as the process of the abolition of private property: “For the abolition of the idea of private property, the idea of communism is quite sufficient. For the abolition of private property in actual reality, however, actual communist action is required.”

One point deserves mention, as we consider it important. It concerns the use of the terms “socialism” and “communism”, whose familiar meanings did not become established at once. In the Manuscripts of 1844 quoted above, “socialism” still refers to a society in which private property is not in the process of being sublated but has already been sublated. Communism, by contrast, is understood as the process (act) of negating (sublating) private property. That is precisely why it is regarded here not as a form of society but merely as the dynamic principle of the immediate future.

This circumstance sheds light on why Marx and Engels raised precisely these questions, and set them out in precisely this way, in the “Manifesto of the Communist Party”, written just three years after the “Manuscripts”. The roots of the Manifesto’s central proposition become more clearly visible: “But modern bourgeois private property is the final and most complete expression of the system of producing and appropriating products that is based on class antagonisms, on the exploitation of the many by the few. In this sense, the theory of the Communists may be summed up in the single sentence: Abolition of private property.”

This sublation, as already noted, presupposes the preservation of all the cultural wealth, active capacities and needs developed on the basis of private property. It is carried out not for its own sake but for the existence of a communist society with its corresponding property and form of intercourse.

Whatever form communist society may take, man as a natural being must live from nature and appropriate nature: everything that man uses and consumes is made from natural material, by means of natural processes put to man’s service, or is itself those natural processes. Twenty years after writing the “Manifesto”, Marx wrote in the Introduction to the “Economic Manuscripts of 1857–1859”: “All production is appropriation of nature by an individual within and through a specific form of society. In this sense it is a tautology to say that property (appropriation) is a precondition of production. But it is ridiculous to leap from this to a definite form of property, e.g. private property (which, moreover, presupposes an antithetical form, non-property, as a condition).”

Communist society corresponds to a common, that is, social form of property in all the social means of production and, accordingly, to common social management of all social production as a single whole. The social character of production, developed on the basis of private property, must be matched by the social character of appropriation of the products of that production. Marx and Engels wrote about this dozens of times in various works. This is perhaps the only proposition that modern communists have remembered, even if they do not understand its significance.

To set this historical movement in motion in the most advanced countries, Marx and Engels proposed the following measures: “1. Abolition of property in land and application of all rents of land to public purposes.

2. A heavy progressive tax. 3. Abolition of the right of inheritance. 4. Confiscation of the property of all emigrants and rebels. 5. Centralisation of credit in the hands of the state, by means of a national bank with state capital and an exclusive monopoly. 6. Centralisation of all transport in the hands of the state. 7. Extension of factories and instruments of production owned by the state; the bringing into cultivation of land and the improvement of the soil in accordance with a common plan. 8. Equal liability of all to work; establishment of industrial armies, especially for agriculture. 9. Combination of agriculture with industry; promotion of the gradual elimination of the distinction between town and country. 10. Public and free upbringing of all children. Abolition of children’s factory labour in its present form.”

As we can see, given the stated aim of the abolition of private property, these measures do not look all that radical; rather, they are general, as befits a manifesto. Manifestos must be clear and comprehensible, and they must indicate the main direction of development. Moreover, the authors themselves maintained that the measures depend on concrete-historical conditions and may therefore differ. We will discuss these measures specifically in one of our upcoming issues, through the prism of the experience of 20th-century socialism.

For now, a few more words need to be said about what private property is, and what exactly communists set out to abolish while preserving all of its positive content, so that “in place of the old bourgeois society, with its classes and class antagonisms,” we shall have “an association, in which the free development of each is the condition for the free development of all.”

Private property is the product of unfree labour and, at the same time, a means of production of the unfree character of labour: labour in which neither the process nor the product of labour belongs directly to the collective of working people; labour that is, in itself, a matter of indifference to the workers and serves merely as a means of achieving external ends (money). The essence of private property lies in this character of labour. It reaches the limits of its development in capitalist society with its developed commodity production, thanks to the social division of labour being carried to its extreme, turning the products of labour into commodities (in commodity production).

For a long time, the social division of labour furthered the development of humanity. It raised the productivity of labour and created the material basis for transferring the machine-like functions that people are forced to perform because of underdeveloped tools to those very tools, to machines. But it does not merely unite people along the lines of “Vasya is a fitter and Petya is a programmer, each does his own job, and above them stands a specialist manager who directs them for the common good.” No participant in the process pursues such a goal. Under this state of affairs, all people end up connected to one another externally—through the products of their labour. They are forced to exchange their product for one they do not have, while remaining indifferent to one another in the labour process. It does not matter who makes this product or how; what matters is that I need this product of someone else’s labour.

The social division of labour has passed through several major stages that played a positive role in the development of humanity. They deserve to be discussed in a separate work devoted specifically to this question. Here it is only important to note that social development is a product of history. Marxists reject outright all prejudices about the biological or natural suitability of members of society for particular kinds of labour. These prejudices do not withstand even superficial scientific criticism. They are at odds not only with what is known about the history of this question, but also with the history unfolding before our very eyes. Every year a multitude of new professions appear and disappear. Parents often do not understand what their children’s work is actually about, and they certainly could not have passed on an aptitude for it through inheritance, along with their genes.

From the moment labour becomes partial and the products of people’s labour become commodities, that is, begin to be produced specifically for exchange, these products begin to dominate people. Things—the products of their own labour—and money, as the representative of all things, acquire a well-known power over living people. And from the moment labour power becomes a commodity, such domination is accordingly acquired by those who have enough of this money to buy means of production and labour power. These people own living human labour. Economically, they are merely representatives of capital. Their function is to attach other people’s labour to capital so that the value of capital increases. Ultimately, this leads to a situation in which almost all members of society, in one way or another, work for the sake of increasing capital—the self-expansion of value. “Only such products can become commodities with regard to each other, as result from different kinds of labour, each kind being carried on independently and for the account of private individuals.” Therefore the abolition of private property is the abolition of commodity production in general, and not only of capitalist commodity production.

Value strives for self-expansion. This is the law that governs a society in which everyone is chained to their own partial activity—in other words, a society of developed commodity production. The laws of modern states are the natural juridical expression of this objective law of production, which is based on unfree and fragmented labour.

Up to a certain point, capital's drive for self-expansion is the driving force of social development. But it calls into being social forces that become cramped within the confines of capital. Using them exclusively for the purpose of increasing capital leads to catastrophic consequences.

These forces must be managed—and indeed can only be managed for the good of society—by everyone together, for the development of all, according to a common rational plan. And for this it is necessary to master universal human capacities (to acquire universal sensuousness, thinking and morality), going beyond the confines of one's profession, which do not allow one to make sense even of a “neighbouring” profession, let alone of social processes as a whole. Therefore the abolition of private property is equivalent to the abolition of the unfree character of labour. It is accomplished as free communist activity develops and becomes the common possession of all people.

To private property, communism opposes common property as a form of people's joint practical relationship to the means of production and natural resources. Common (Social) Property ensures the most efficient use of productive capacities not only for society as a whole but also for each of its members. It presupposes active participation and an active will with respect to social resources on the part of every member of society. Accordingly, if private property is a product of the division of labour (its fragmentation into ever more particular kinds of labour)\footnote{It was precisely because the social division of labour in the USSR had not been overcome (the tendency towards overcoming it had not become predominant), and the proletariat was reproduced as a proletariat, that is, as a class of bourgeois society, that the complete restoration of the rule of private property took place.}, and is in fact this very division, then common property presupposes, as noted above, the socialisation of labour and the elimination of its division. Social property presupposes free human activity by the standards of truth, goodness and beauty. It is this that constitutes the real economic basis for a real, and not merely formal, socialisation, and for a real coincidence of the private and social interests of all members of society.

Common or social property should not be understood as the property of collective farms and cooperatives, or of other groups of people who produce commodities and trade them with the rest of the world. These do not go beyond the bounds of private property, since they do not go beyond the bounds of commodity production. The members of such enterprises are connected with society as commodity owners and commodity producers. The very nature of commodity production also determines the relations within these collectives, since it puts the production of value first and makes profit the goal of the collective's work.

Nor should the germ of common property be sought, for example, in the property of the family. In the 19th century, some socialists imagined that social property could arise through the “expansion” of family property to the scale of society as a whole. Society was pictured as a large family with the methods of distribution characteristic of a family. But the family is the embodiment of the private character of appropriation of the products of labour. In its modern form, it is itself a product of the development of private property. Insofar as the family carries out the production and reproduction of human beings as labour power (the human being as an economic agent of a society of commodity production), it is wholly embedded in, and entirely reproduced according to, the laws of motion of private property. Insofar as it reproduces all its members and serves as the basic consumer “unit”, it is wholly and completely determined by capitalist production. Economically, this is presupposed both in the form of wages and in the form of profit. Moreover, the modern family as an economic unit not only does not exclude but directly presupposes the exploitation of man by man, which is incompatible with communism. Therefore, family property cannot be regarded as the basis of social property—the property characteristic of communism.

And still less can communist property, which might serve as a nodal point in the movement towards communism, be found in all sorts of isolated communities unconnected with social production (for example, religious communes based on a natural economy).

But where, then, can we see the common property characteristic of communism? Clearly, only where there is free human activity in its directly collective form. Such activity, viewed from the standpoint of the use and disposal of resources, from the standpoint of appropriation by means of a social form, is precisely the actual existence of common property.

In the previous chapter we gave only a few examples of such communist activity in socialist and capitalist society. If we consider these and other examples from the standpoint of the exercise of will in relation to the subject of labour, the means of labour and the product, we are dealing with forms of social property and the social character of appropriation.

For example, a publicly accessible database, by its very nature, can no longer be anything other than social property, to say nothing of a common system of economic management. The same applies to free software and to the various services that presuppose joint production and use. There are also resource-sharing servers that would be inefficient and unprofitable to operate on the basis of private property—for example, shared libraries and repositories for films and music, which presuppose that participants both receive and distribute products on these services (although the example of software and databases is better, since it is more a matter of production than of consumption).

It is clear from the above that social property arises from the real movement of production even where there is as yet no formal socialisation of the means of production, under the conditions of capitalism. Formal socialisation—the seizure of political power by the proletariat and its taking possession of the means of production—will open up enormous scope for the development of real social property. Note that wherever there is real, and not merely formal, social property, the ground disappears for attempts to appropriate a portion of social resources for oneself, or, conversely, to squander them and treat them carelessly as though they belonged to “no one”.

Under formal socialisation, some external regulation is required—at least in the form of moral consciousness (in reality, in the form of law and a system of coercion that ensures its enforcement)—in relation to the means, subject and products of labour, so as to rule out such things. Under real socialisation, the need for such regulation falls away insofar as a culture of relating to the means of production becomes a simple condition of their use.
\ifSubfilesClassLoaded{\setcounter{page}{36}}{}
\section{Necessary Labour and the Conditions for Communist Activity}
\label{sec:06_necessary_labour}

In communist society, human life is not mediated by private property and the exchange of commodities; consequently, it is free from everything required to secure the conditions for the existence of private property. This gives society considerable advantages and saves its labour power. In the “Speeches in Elberfeld”, Engels offered a number of examples that show most vividly how much human labour can be saved in a society where the interests of its members are not opposed to one another and the products of activity do not turn into commodities confronting one another on the market.

The labour needed to keep the market running, the market through which social resources are distributed, is replaced by more efficient direct distribution. This makes it possible to stop performing actions that are senseless from the standpoint of organising production on a social scale, and to stop squandering people’s lifetimes on them. Producing packaging that is unnecessary for simply storing goods, carrying out shipments tied to sales rather than to delivery to the end consumer, producing goods that will later turn out to be surplus on the market and go unsold, and, still more, producing things that harm people yet bring enormous profits in a society of commodity production (drugs and the like)—none of this makes any sense in a society where production does not have a commodity character.

Abolishing the private character of appropriation also yields colossal savings of living human labour. This is the same as in the case of commodity production, but viewed from the side of the labour spent bringing products to final consumption. The examples of everyday life and transport, which we have already written about, positively cry out that it is possible to free up a vast amount of human labour that is currently expended in an inefficient, individual way.

Modern capitalism demands human effort aimed at sustaining the capacity to produce surplus value. There is more and more labour of this kind. The production of surplus value has, of necessity, been supplemented by the production of consumption and of a way of life geared towards consuming particular commodities. This alone shows the growing (and, compared with the times of Marx and Engels, greatly increased) possibilities for liberating people, who will be able, on the basis of collective property, to engage in producing what society needs for its development.

Social time is spent on maintaining the restrictions that private property imposes on the life of society and on eliminating its consequences that are harmful to society, including crime: the police, the standing army, officials, and all manner of guards of private property. All this labour power, too, can be freed up once private property in the means of production is overcome, and redirected to performing labour in which society has a direct interest.

All of this in itself, even on the basis of already existing technologies, will reduce socially necessary labour time and create room for free time to exist for genuine human life—that is, for communism proper.

Another important tendency also points to movement in the direction of communism. It is maturing within capitalism and is held back by it only because capital is the dominant social relation. “The great historic side of capital is the creation of this surplus labour, superfluous from the standpoint of mere use value, of mere subsistence of the worker, and the historic destiny of capital will be fulfilled when, on the one hand, needs have been developed to such an extent that surplus labour itself, labour beyond what is absolutely necessary for life, has become a general need arising out of people’s individual needs themselves; and when, on the other hand, general industriousness, thanks to the strict discipline of capital that successive generations have passed through, has developed as the general property of the new generation; and when, finally, this general industriousness, thanks to the development of the productive forces of labour, constantly spurred on by capital with its boundless passion for enrichment and acting under the only conditions in which it can realise this passion, has brought it about that, on the one hand, the possession and preservation of general wealth require only a comparatively small amount of labour time from society as a whole, and, on the other hand, the working society relates scientifically to the process of its progressive reproduction, its reproduction in ever-growing abundance; hence, when the kind of labour comes to an end in which a human being himself does what he can make things do for him, for man,” wrote Marx. We would add that this would mark the completion of the entire period of the existence and development of the social division of labour.

The idea of the automated factory, which emerged as early as Marx’s time, and the idea of making wide use of nature’s potential, through knowledge of it, for the good of all people, capture the really existing tendencies towards creating the conditions for communism. Communism will advance as things increasingly do the work they can do for human beings and without them, as there is less and less of the unfree human labour necessary for reproducing human life, and as the government of persons is increasingly transformed into the administration of things. Every person will be able to devote more and more time to the free play of his bodily and mental powers or, in other words, to doing what is “after his own heart”, to the benefit of society as a whole. This is precisely what constitutes the essence of communist activity.

The historical role of capital as a social relation essential to the development of the productive forces cannot be fulfilled under capitalism. Under the undivided rule of capital, capital constantly reproduces the thing-like character of labour, time and again throwing the development of the productive forces back wherever a tendency to overcome that character has emerged. For this reason, under capitalism science cannot fully become a direct force of production, and as the thing-like character of labour is reproduced, so too, among others, are its most backward forms. The complete abolition of capital as a social relation is the most important question of the period of transition to communist society. In this period, the movement of capital towards handing purely material production over to machines must be consciously carried through to the point where human activity is taken beyond the bounds of the thing-like character of labour.

Here, however, we must not forget that capital at the same time also reproduces the opposite tendency, one that leads society back to capitalism. It is this tendency, inseparably bound up with the first though it is, that socialist society can and must openly combat if the transition to communism is to take place at all. The front line between the old capitalist and the new communist objectively runs between those spheres where all the material preconditions exist for handing thing-like labour over to things once and for all, and those where this cannot yet be done without first carrying the division of labour to its utmost extreme and thereby creating the physical preconditions for transferring it to machines. Wherever the complete automation of the purely material production of things is possible, capital as a social relation, together with the corresponding character of labour, loses its relevance and therefore objectively can and must be abolished. And until this is done, capital will continue to reproduce itself as an essential social relation.

The notorious question of the anti-communist philistine, “Who will clean the toilets if everyone gets to do what they love?”, loses all meaning. Its rhetorical force is deflated not only because there is a vast difference between occasionally setting aside one’s favourite pursuit to perform this unpleasant but necessary social duty and being chained to it for life, but also because self-cleaning toilets already exist. It is only the dominance of private property, and the fact that under these conditions the labour power of cleaners, that is, the lifetime of human beings, turns out to be cheaper than things, that prevents their widespread adoption.

Under communism, by contrast, the human being, that is, the time and fullness of a person’s life, becomes the highest value, and all things matter only insofar as they secure a genuinely human life for every member of society. The very fact that a living person must spend their lifetime on uncreative labour that could be handed over to machines is intolerable for a society in which the human being is the highest value. This state of affairs must be abolished, and only then will it be possible to speak of communism developed on its own foundation.

It must be acknowledged, however, that during the period of transition to communism, necessary labour that lacks full freedom will persist for some time, until the appropriate technologies for its transcendence have been developed and introduced everywhere. But eliminating the rule of capital will make it possible to limit the working day to necessary labour, excluding labour that has meaning only within the movement of value towards its self-expansion, the results of which are time and again destroyed during crises or simply written off as “costs”. At first, necessary labour (in the sense of labour that is not yet fully free) will expand its bounds. This is needed to provide social funds and the conditions for further development. But the general direction of society’s movement towards communism, which includes the corresponding development of the productive forces, presupposes reducing such labour to ever smaller proportions. Beyond a certain point, such labour can be left out of account altogether, since it no longer has any serious influence on the general course of human life. The most important thing about necessary labour in the transition to communism, as discussed here, is that this not yet entirely free labour is labour that produces free time and a human being capable of acting freely on a social scale.
\ifSubfilesClassLoaded{\setcounter{page}{40}}{}
\section{Communism and the Management of Social Life}
\label{sec:07_communism_management}

The production and distribution of free time during the period of transition to communism require the state. However, the state is necessary only insofar as the production of free time coincides with the production and distribution of the things people need, and insofar as this production of things is not yet carried out in a communist way. Free time itself, emerging as genuinely free (the time of communism), escapes any state management external to it.

Within the space of free time, the state does not exist even now, under capitalism. Wherever free communist activity takes place, management is carried out by all the members of a free association of people. As the share of communist activity grows, the state withers away.

Promoting the production of free time and of directly collective forms of free activity is an act of abolishing the conditions for the existence of classes during the period of transition, and to the same extent an act of the self-abolition of the state. The state is merely an apparatus for securing the interests of the ruling class, exercising management in the interests of that class, and this primary function falls away as classes, including the ruling class, are abolished.

The task of the state during the period of transition thus comes down to ensuring the restriction (abolition) of the commodity character of production and to concern for creating the conditions for the free communist activity of the members of society.

The first of these tasks confronts capitalist states as well. Capitalism constantly breeds wars, and the logic of war demands a system of direct distribution of products intended for personal consumption. A state that fails to cope with this task is doomed to destruction (either on its own, as a state, or together with its population). The task of establishing such a system entails, at a number of points, the curbing of individual capitalists in the interests of capitalists as a class, which cannot always be achieved in full or with a sufficient degree of effectiveness.

In war, personal consumption—of both food and clothing supplies—cannot be provided for through the capitalist private character of appropriation of the products of labour. A soldier, even a mercenary, must be supplied by his employer during wartime with, at the very least, weapons, food and uniforms. Otherwise, the army degenerates into gangs of brigands. But such gangs are incapable of waging modern warfare. To wage war, a modern army needs high-tech weaponry, the means of personal consumption necessary for reproducing the living personality of the soldier, and personal protective equipment. And even if all of this is secured through the outright plunder of the population of captured territories, it cannot be the personal affair of individual soldiers. It is merely a condition that enables the soldier to do his job. But this condition differs from the conditions of wage labour, which is based on the private character of appropriation. In classical capitalist relations of the purchase and sale of labour power, the appropriation of the products of labour, in all of its aspects that are essential and significant for the social reproduction of the individual, is placed outside the sphere of production based on wage labour. Personal consumption is determined by the private character of appropriation of the products of labour, mediated by money. And this, in turn, guarantees that the worker will be compelled, again and again, to return to work in order to provide himself with everything he needs. In war, things are radically different. Here, consumption cannot be left to individuals themselves, even formally.

Thus, modern warfare is, among other things, a competition between forms of non-commodity production that took shape under capitalism itself—the direct production of man. This applies both to the army as a military force proper and to the labour army, which is formed in various ways. The emphasis shifts. Whereas in the normal course of capitalism what comes to the fore is that a person must do his work in order to live, here, on the contrary, a person must live in order to do his work. The state is forced to take responsibility for this. Moreover, commodity-money relations become an obvious obstacle to the performance of these functions.

War makes the old, pre-war ways and forms of human activity impossible. Under such conditions, the state objectively grows stronger, takes on the functions of distribution in the process of the direct production of man, both in the army and on the home front, and works to establish and maintain economic ties. But the state, as an apparatus of suppression of one class by another—or, more precisely, of the majority by a minority—is in principle incapable of fully accomplishing these tasks, let alone setting itself tasks that do not serve the interests of the ruling class. Yet it is precisely tasks of this kind (those at odds with the interests of the ruling class) that multiply for society during war. They objectively confront people because their very physical existence depends on solving them. The state often does not even attempt to solve these problems, leaving them to a wide variety of non-state organisations and associations of people. This compels millions of people to seek new forms of organisation that would take over state functions. And these bodies of social self-government, which take on the functions of the state, strictly speaking, in their tendency both are and are not a state.

The most important and formative of the problems whose solution requires new organisational forms in opposition to the capitalist state is the prevention and ending of war, along with minimising the losses and destruction it inflicts on society. Accomplishing this crucial task runs directly counter to the interests of capital, which has not yet extracted from the war all the benefits it could. Let us stress that what may carry historical meaning here is not war as such, nor violence as such, but resistance to the violence of the reactionary class, violence carried out by factions of that class in the name of their own selfish interests, and by the hands of another class at that. It is in resistance and in the forms of resistance (though not only there), rather than in war as such, that the most important tendency towards self-organisation on the approaches to revolution lies. Without such self-organisation, the methods of the direct, unmediated by things and money, production of man that have been developed within capitalism, along with the economic mechanisms for implementing them, cannot become the starting point for society’s transition to new foundations of development.

These forms of self-organisation for managing public affairs, generated by the tendency towards war and by war itself (as opposition to war), are transitional and can move in three directions: 1) integration with the state; 2) degeneration and disappearance (integration does not always mean degeneration); 3) taking over all the functions of the state and turning into a new state power. These organisational forms are thus, objectively, also forms of class struggle wherever they are forced to oppose the state by taking over the functions of particular state bodies, and hence stripping the apparatus of the ruling class’s political domination of its powers. And while under certain conditions this actually comes as a relief to those bodies, since they themselves are unable to perform their functions, any further movement along this path inevitably runs up against the opposition of the state and its attempts to restrict the functions and influence of the population’s forms of self-organisation. The outcome and character of this confrontation will depend on the specific circumstances of the struggle and the actions of the contending parties.

Looking at the history of the twentieth century, one of the most effective forms of self-organisation to take on the functions of the state were the Soviets. The Soviets of Workers’, Soldiers’ and Peasants’ Deputies served simultaneously as representative-legislative, executive and supervisory bodies. And although the Soviets first appeared as early as 1905, their real development into a state force, dual power, let alone the slogan “All power to the Soviets” and the concrete steps from that slogan to its implementation, could only arise amid the political bankruptcy of the old state in the course of an exhausting war. After the February Revolution, the bankruptcy of the state as an apparatus claiming to organise the life of society became obvious, since not only the tsarist government but also the Provisional Government proved powerless in this respect. At the same time, the Soviets were rising and actively involving themselves in the affairs of both local self-government and state administration. A situation of dual power was taking shape both on the home front and in the army. For this reason, it is especially important to bear in mind that the situation of spring and early summer 1917 could arise again in the future. At that time, a peaceful transfer of power to the Soviets was possible as a step along the path towards social revolution, that is, towards creating the conditions for socialism. This possibility existed thanks to two circumstances: the strength and social significance of the Soviets themselves, and the fact that the population had been armed as a result of the war. The slogan “All power to the Soviets” was, at the time, a statement that such a peaceful path of revolution existed. Only after the July Days of 1917 did the peaceful path become impossible. But the very fact that such a possibility existed is what matters most to us. After all, if such conditions arise, it is vital not to miss them, especially since, given the development of modern weaponry, the peaceful path may well prove to be the only possible one for revolution.

Under present-day conditions, other forms can emerge, and are emerging, which are by no means always connected with war but are always connected with the need to solve important social problems. The essence of these transitional forms is the same everywhere: the population directly takes upon itself the functions of state power, thereby bringing about the transcendence of that power in the sense customary under capitalism.

Let us now turn to the management of social life under communism. In general, the management of social life comes down to planning and adjusting the plan in light of real changes, the expression of will (issuing directives), accounting and control over the managed process, and direct intervention where necessary.

Planning under communism is planning for the development of the human being, who becomes the principal aim of social production; it is planning for the production of new relations between people. Such planning coincides with forecasting, in the sense that the most reliable forecast is the one towards whose fulfilment systematic effort is directed. Society must learn to perform all these functions without an organ of “management” standing above it, which is in reality an organ of oppression, and which is called the state.

In the course of the becoming of the new collectivity, however, all the functions of the state that relate to the production of things must be handed over to things. In the period of transition to communism, the state does not simply wither away; it withers away only if it works towards its own withering away, to the extent that it takes care that people can do without it. All the remnants of unfreedom, including the ways of utilising time with which the existence of the state is bound up, stand opposed to free time as the space of human development, and stand opposed to free human activity. Where complete freedom is achievable, there is no place for the state. But where complete freedom is not yet directly possible, the existence of a certain type of state not only registers the fact of this impossibility but also serves as an instrument for overcoming it.

The same argument is advanced against the withering away of the state as against the sublation of the division of labour. Its essence is that social life is becoming very complex and is growing ever more complex, and that managing such a complex life therefore requires a differentiation of functions and a special apparatus. But it by no means follows from this complexity that special detachments of officials and armed men, set above society, are necessary. Were it not for the split of society into classes, a self-acting and self-governing association of people would be possible, for all its complexity corresponding to the level of complexity of social life, as Lenin, following Marx and Engels, insisted.

All the more so since, when we speak of communism, we mean a society in which the division of labour has been overcome, in which people are more or less equally developed in terms of health, sensuousness, thinking and morality, and have a polytechnical, scientific-theoretical, universal education at the level necessary for society today (that is, the highest), and can therefore master any particular “profession” in a short time and apply their efforts wherever they are needed and wherever they themselves find it interesting, even if one person devotes most of his life to studying birds and another to ballet.

But it is precisely for this reason that under communism everyone takes part in the management of social development, using the technologies achieved and relying on the culture of thinking that humanity has developed. This is the most important condition for the universal development of all members of society: throughout their lives, everyone must be able to manage, and be in the habit of managing, social affairs—all the affairs in which they themselves take part. Such a society is more realistic than the fantastical class peace and “harmonious relations” between classes propagandised by today’s bourgeois ideologists. On the contrary, Marxism insists that the very existence of the state shows that such peace is impossible as long as classes exist.

The withering away of the state requires the appropriate social and technical preconditions. As for the technical preconditions, with good universal polytechnical education grounded in modern technology, organising the management of social processes would not be all that difficult in a society without private property (which is why the formal socialisation of the means of production is the most important precondition). Of course, social barriers must be removed for this, which likewise demands serious effort from society and cannot be achieved overnight. The main point is that, purely technically, it no longer seems impossible even at the present stage. All information on the state of production and distribution in every sphere of social life in real time, with the possibility of tracing chains back to their very beginning, accounting of resources, needs and technological capabilities, the possibilities of communication and voting, and so on and so forth, from the lowest level (direct intervention) up to the highest level of collective decision-making, both in local communities and in society as a whole—all of it can be available in real time to every member of the association.

Voting and the other familiar expressions of democracy are needed only at first. With the transition from the government of persons to the administration of things and processes, the need for voting gradually falls away, giving way to the objective necessity inherent in things and in the interactions between them. Lenin pointed this out in The State and Revolution in his detailed examination of the legacy of Marx and Engels: “In speaking of the ‘withering away’, and the even more graphic and colourful ‘dying down of itself’, Engels refers quite clearly and definitely to the period after ‘the taking possession of the means of production in the name of society’, that is, after the socialist revolution. We all know that the political form of the ‘state’ at that time is the most complete democracy. But it never enters the head of any of the opportunists, who shamelessly distort Marxism, that Engels is consequently speaking here of democracy ‘dying down of itself’, or ‘withering away’. This seems very strange at first sight. But it is ‘incomprehensible’ only to those who have not thought about democracy also being a state and, consequently, also disappearing when the state disappears. Revolution alone can ‘abolish’ the bourgeois state. The state in general, i.e., the most complete democracy, can only ‘wither away’.”

In today’s era of rapid technological development, organising the withering away of the state does not seem such a difficult matter from a purely technical standpoint. Under such conditions, accounting and control in the administration of things and processes become a matter of technology, and their results become accessible to everyone. Here the need for hierarchy falls away, and the precondition for these functions to become separated into a distinct profession disappears. With the introduction of automated management systems, and with full access to databases and the ability to work with them guaranteed for everyone, the need disappears for a “special army of officials” who could stand, even “in appearance”, “above” society as a special “apparatus”.

In this way, one of the most important tasks of the transition to socialism will be accomplished: establishing accounting and control, the importance of which Lenin constantly stressed: “Accounting and control—that is mainly what is needed for the ‘smooth working’, for the proper functioning, of the first phase of communist society.” What we are talking about is accounting and control over social production of a kind whose very organisation already means the organisation of new forms of collectivity: “If really all take part in the administration of the state, capitalism cannot retain its hold. The development of capitalism, in turn, creates the preconditions that enable really ‘all’ to take part in the administration of the state. Some of these preconditions are: universal literacy, which has already been achieved in a number of the most advanced capitalist countries, then the ‘training and disciplining’ of millions of workers by the huge, complex, socialised apparatus of the postal service, railways, big factories, large-scale commerce, banking, etc., etc.”

We, for our part, can emphasise once again that these preconditions now also include the development of computer technology and the potential it holds, which in principle cannot be realised under capitalism, where there is no single social economy and the connection between producers is mediated by the market. This potential for organising the management of social life was explored in the work of Victor Mikhailovich Glushkov, who was developing the nationwide system of economic management in the USSR. Glushkov considered the management of social production as a whole by means of control systems in which everyone has access to a highly detailed, real-time picture of production, while wherever human decisions are not required, the system acts automatically according to schemes set in advance by people in line with social goals.

The more technology develops, the easier it will be to organise management and to ensure that all members of society take part in it. As for the substantive side, it requires a corresponding development of the members of society. The difficulty is that the movement towards communism has to begin not with the people of the communist tomorrow, but with people shaped under the conditions of the division of labour, who overcome their limitations through education and collective activity.

Polytechnical education becomes a necessity for the members of a free association. The idea of polytechnism has been worked out quite thoroughly, so we can safely refer the reader to the relevant literature. Here we will only note that in order to manage social processes, members of society must possess an integral scientific picture of the world and a universal ability to apply their knowledge. A modern official has absolutely no need for a scientific picture of the world, even in rough outline, to perform his function, and so it simply disappears of its own accord as something superfluous, even if he once had one. An official is merely a living, partial component of an obsolete state machine. “The society that will organise production anew on the basis of a free and equal association of the producers will put the whole state machine where it will then belong: into the museum of antiquities, by the side of the spinning wheel and the bronze axe,” said Engels.

Of course, the very task of managing social processes in such an association, and of consciously building social life, will demand a great deal of people, including well-developed dialectical thinking, well-developed sensuousness, morality, and an adequate level of literacy. But these are not unattainable heights accessible only to a chosen few; they are something that, under the right conditions, will become the possession of everyone. As people are drawn into active social life as its subjects, they will come to the forefront of human culture. This is how we understand Lenin’s famous aphorism that every cook must learn to govern the state. The only obstacle to this is the mode of activity of people today, which turns one person into a state official and another into a cook.

Communist society is a free association of people governing its own development according to a rational plan.

How exactly will it govern itself, and by what specific means and organisational measures? We cannot answer this question. That is a matter for the association itself, which will of necessity have to develop the appropriate organisational forms under concrete historical conditions. Self-government will develop along with the development of communist society.
\ifSubfilesClassLoaded{\setcounter{page}{47}}{}
\section{Communism and the Personality}
\label{sec:08_communism_personality}

The problem of the personality is one of the central problems of communism. Under communism, the significance of the personality rises to heights never before seen. “In place of the old bourgeois society, with its classes and class antagonisms, we shall have an association, in which the free development of each is the condition for the free development of all,” wrote Marx and Engels in the “Manifesto of the Communist Party”.

All of preceding history must be understood not simply as a particular sequence of events, but as the process of the becoming of the human essence. It then becomes clear that this is, at the same time, the process of the becoming of a particular type of personality, one that corresponds to the human essence. What should this personality be like? This question will be resolved each time practically and in concrete-historical terms, but in any case it is a personality that consciously creates history within a collective of other creative personalities like itself. In a free association of people, their entire life depends on creative activity, which is the highest expression of the personality that carries it out.

The personality is a social phenomenon from beginning to end. An individual is not born a personality but becomes one, entering into social relations as an actively engaged subject. Initially, human life is given to the newborn from outside—from the social space in which they live, with all the variations of the relations of human being to human being and of human being to nature. Only by assimilating, that is, by making their own, the forms of activity—or, what amounts to the same thing, the forms of culture, which are in no way inherent in the individual from the start—do they become a personality. Everyone who comes into the world finds these forms already in place. For the biological organism they are external, but only for the organism, not for the personality. The world of activity, as the world of human culture, is the real world of the personality. In communist society, everyone must make this world their own, and that means learning to change it.

The communist personality will, in its practice, actively create social relations rather than simply being their passive product. The problem is to create conditions in which the personality, in its individual development, would travel in the shortest possible time the path of society’s becoming and arrive at the forefront of social life. Each time, one must seek and find concrete forms of collectivity that would bind the personality to society in such a way that it could not help being an active member of it and living the full life of society. Since society is constantly developing, such forms of “joint-separate activity” cannot be found once and for all. The management of social development under communism consists mainly in learning to change, in good time, the forms of connection between the human being and society so that they compel people to act humanly in any new concrete-historical conditions.

Communists are accused of wanting to standardise the personality, to make people “all the same”. In making this accusation, the supporters of capitalism are in fact giving away the secret of modern capitalist society. Under capitalism, individuals by and large turn out so standardised that their individuality is defined by mere trifles. This bad individuality has nothing to do with true individuality, understood as the individual expression of the universal human culture a person has assimilated. Capitalism draws people into standard practices of production and consumption and standardises their behaviour both at work and outside it. People are already “all the same”. Communism, by contrast, presupposes the creation of a social space for the development of individuality.

Communists see the need to create a society of universally developed creative personalities. In such a society, the individuality of each person will be the active expression and embodiment of social development. In what sense, then, can we speak of standards here?

Clearly, a person cannot be a universal creator without having assimilated the minimum of human culture necessary for this, and so we must speak of a minimum of conditions that society must provide to each of its members so that they can be a universal creator; there must be a minimum standard for these conditions. What does this minimum consist of?

It includes: a certain level of development of the healthcare system, one aimed specifically at preserving health and creating the conditions for a healthy life for every person, rather than merely treating illness; a quality and level of education that is compulsory for everyone; a system of public catering and, more broadly, of all publicly accessible services capable of abolishing everyday life; and provision for childhood and old age. In one of his works, Lenin wrote that socialism begins with a glass of milk for every child. Obviously, the presence of that glass does not in itself mean that society is developing in the direction of communism, but without children being well fed, it certainly cannot develop in that direction. Nor can it develop without everything listed above (the list, of course, is not complete. It should be drawn up by people with expertise in various fields, and it will expand as society develops).

This minimum of humane treatment of one person by another cannot be fixed once and for all. It changes depending on the conditions in which the transition takes place, that is, on the level of development of human culture. Today’s minimum, the one that would give society the impetus to develop towards communism, must necessarily include higher education encompassing not only computer literacy at the level of programming but also the scientific foundations for understanding all production processes—polytechnism (for everyone!). Ensuring access to all modern means of communication is likewise essential.

It is important to note that even the most developed capitalist countries cannot boast of universal literacy among their entire population or of universally accessible medical care.

The necessary minimum of humaneness in how people are treated is, as already noted, a historical matter. Yet this minimum must also include the need for its own constant expansion. For a communist society, in which the free development of each is the condition for the free development of all, to exist and develop, this minimum must continually incorporate ever newer, most advanced achievements of humankind. Communist society will consist of intelligent, healthy and talented people, who, in the final analysis, are the only ones capable of organising the social reproduction of their lives on rational foundations.
\ifSubfilesClassLoaded{\setcounter{page}{50}}{}
\section{Communism and the Family}
\label{sec:09_communism_family}

“According to the materialist conception, the determining factor in history is, in the last instance, the production and reproduction of immediate life. This, again, is of a twofold character. On the one side, the production of the means of subsistence: of articles of food, clothing and shelter, and the tools necessary for that production; on the other side, the production of human beings themselves, the propagation of the species. The social organisation under which the people of a particular historical epoch and a particular country live is determined by both kinds of production: by the stage of development of labour on the one hand, and of the family on the other,” wrote Frederick Engels in The Origin of the Family, Private Property and the State.

In the previous chapters, we discussed how, under communism, the production of the necessities of life is automated and subordinated to the production of human life in all its fullness. Society consciously creates its own development and its own forms of collectivity, expanding the content of free activity. The family, as an economic unit based on the private character of appropriation of the products of labour, loses its significance and disappears during the period of transition. “Private housekeeping is transformed into a social industry,” wrote Engels. We saw above how much social time this saves, freeing people from the burdens of everyday life. However, socialisation concerns more than just consumption.

The main revolution in family relations lies in the social upbringing of children. This revolution has long been brewing, and movement in this direction already exists under capitalism. Housekeeping has begun to be “socialised”, and everyday life has long since become, in many respects, a branch of the social economy. We cook mostly from semi-prepared foods, buy our clothes in shops and wash them at the laundry, even though not long ago all this was women’s work done at home. The modern man, too, no longer needs to carry water, chop firewood or build the house his family lives in. All of this has long been a high-tech social undertaking.

But the upbringing of children still refuses to become a sphere of production, even though public institutions for this purpose appeared long ago and have been actively developing for at least two centuries. And the better-off have always sent their children precisely there and hired private tutors, while they themselves merely enjoyed the comforts of family life. So it was in the Soviet Union, too: one of the first achievements that socialist working people spoke of with such pride was the organisation of factory nurseries and kindergartens. In other words, the need to entrust the upbringing of children to those who know how to do it, and to win free time for what one considers truly necessary and important (for some, incidentally, that may be the upbringing of children—why not?), is something rich and poor alike share, with only a few differences of a petty-bourgeois, sentimental nature. Capitalism has long been destroying the traditional family, and this process has no moral roots, only economic ones. At present, the promotion of family values is driven, on the one hand, by the achievements of socialist revolutions, which banned the exploitation of child labour, and on the other, by capital’s need for a market in children’s commodities and education. But there is no real historical or social necessity to preserve the family in its current form as an economic unit.

Thus, as society transitions to communist relations, the sphere of the upbringing of children will become a sphere of social production. Upbringing becomes genuine social creativity—free activity by the standards of truth, goodness and beauty. This increases many times over the responsibility of educators both to children and to society as a whole. As A. S. Makarenko once wrote: “We must have no childhood catastrophes, no failures, no percentage of rejects, even if expressed in hundredths of a unit!” In a number of works, he compared upbringing with other branches and showed how much more dangerous failure is here than anywhere else. Speaking of the family, Anton Semyonovich observed that the task of parents is to raise worthy members of society; they are accountable to society for the “quality of the material supplied”, and in return they receive the warmth of family relationships. But by no means all parents know how to raise children. Moreover, even now, by no means all children have parents or simply a complete family. Herein lies the need for a conscious, high-quality socialisation of upbringing and for the creation of other forms of “family”.

The communist “family” is a collective form in which everyone takes part in producing social life on a correspondingly social scale. It follows that the size of such a “family” may vary, and it certainly cannot be confined within its present-day bounds. Its size depends on many concrete-historical conditions of social production, so various forms will be possible, both within a single epoch in the development of communist society and over the course of its history. Since such a “family”, as the primary collective of the human being in the making, will be founded on the social character of activity and directly determined by it, the principle on which such collectives are formed will likewise be directly social—qualitatively different from what it is now. What is at stake is the abolition of modern marriage as a form, especially since sexual relations and sexual love, on the one hand, and marital-economic relations, on the other, are far from being one and the same thing. That is why one of A. S. Makarenko’s most important and enduring practical discoveries was that the new “family” must draw every person directly into social production as an active member of the social labour collective, rather than as a consumer of social goods or as an “object” of upbringing.

Tracing the rapid, at times barely perceptible changes in this sphere during the very first years of the Great October Socialist Revolution, Alexandra Kollontai, in her work “Make Way for Winged Eros!”, sought to discern and predict the type of personal relations, or more precisely relations of sexual love, that would take hold once social property had completely triumphed over private property, as new social institutions developed and as elements of an infrastructure capable of relieving the burden of motherhood were introduced, increasingly turning it into a branch of social production.

Wherever the new everyday life was established (and not only in the USSR), it truly emancipated women on a scale unimaginable to the rest of the world, giving them time not only for love itself and for “courtship”, but also for rethinking love. These new, already partly communist relations are also depicted in the fiction of the industrialisation era, both in the prose of Alexandra Kollontai herself (“Vasilisa Malygina”, “A Great Love”) and, for example, in Ilya Ehrenburg’s novel “The Second Day”.

Establishing the commune without the registration of marriage (in the 1920s and 1930s, unregistered marriage became the norm, while registration, on the contrary, came to be seen as a sign of mistrust between partners) where the patriarchal family had previously reigned was, of course, getting far ahead of events. But it was a revolution in how lovers, those in love, married people and divorcees perceived themselves. Not marriage as a life for two and for the sake of two, but family life directly shaped by the goals of a rapidly growing and developing society and of one’s immediate (labour) collective—that is the real agenda of love, one that expresses the tendency towards communism.

In capitalist society, the family, as a zone of love and of the responsibility bound up with it, is a “high-risk zone”. Within the walls of their home, spouses are doomed to live through everything that is programmed by the relations of private property, by questions of means, earnings, spending and the rest of life’s prose. But communism is precisely the elimination of such problems.

The dynamics of the withering away of both the old family and the state as such are complex. However, the half-measures of socialism are powerless here: students and young specialists, leaving the commune-hostels, still ended up in flats, in families, within the old frameworks, only now with new social security. The motivation to outgrow the old cells of the “honeycomb” was lost in a society where commodity-money relations were once again winning back their dominant position. The leap towards the commune was followed by a natural, and seemingly also temporary, return to patriarchy. Once again, it was the care of the parents alone, rather than of the whole collective or society, that came to the fore in upbringing. And yet what was emerging was not only “sons of the regiment” (the war had its own say in this), but also children of the communes.

It was in these “children of the communes” that one of the most important aspects of the becoming of communism was embodied. In the transition to communism, children must cease to be divided into one’s own and other people’s. There are no other people’s children (that is, children who are someone else’s responsibility to care for, not mine) in a vast collective where the problems of the “meter” measuring one’s personal labour contribution have been removed.

Finding themselves from the very start in a collective with many men and women, children, while still distinguishing their father and mother, will not see the “edge and end of the world” at the point where mum and dad end and the “neighbours” begin. Through the socialisation that comes naturally in a commune, through the everyday experience of the collective’s care, children will stop “finishing their games in the sandbox of private property”. Moving away from the old form of marriage will make it possible to avoid a number of vestiges reproduced in the psyche and in the consciousness of fathers and mothers, and not only of children. In a large and close-knit collective, a change of partners ceases to be something traumatic if there is no longer love between the members of a couple. In general, in the absence of social economic barriers, everything strictly personal (including the intimate) can only then become truly personal, rather than opposed to the social, once freed from economic, legal and moral shackles.

The new “family” must itself become a directly collective form of communist activity aimed at creating the new man, one that directly links all its members, including children, with society as active subjects.
\ifSubfilesClassLoaded{\setcounter{page}{54}}{}
\section{Communism and Social Consciousness}
\label{sec:10_communism_consciousness}

From the standpoint of Marxism, in carrying out material production, society produces its own consciousness. Humanity has known various historically evolved forms of social consciousness. From mythological consciousness, in which the ideal is not yet separated, or, at various stages, not fully separated, from immediate material activity, philosophy and religion emerge with the development of the division of labour and the alienation of man from his own activity, followed by art and science. Alongside them exist morality and law, which emerged when society came to need external regulation of relations between people with opposing interests. Social consciousness is not the sum of the consciousnesses of separate individuals but the consciousness of society as a whole, separated from immediate human activity: its ideal dimension. Conversely, individual consciousness is that same social consciousness, acquired by a person according to his place and role in social life.

Social consciousness and its forms possess an objective existence in relation to the members of society. To live in human society, people are compelled to reckon with the demands and precepts of morality, law, religion, science, and even art and philosophy, just as they must reckon with surrounding nature and the level of technological development. What is reflected and expressed in social consciousness is not something arbitrary but the mode of interaction between man and nature and between man and man. Social consciousness and all its forms are not an immediate product of nature but a product of its assimilation by human society: a reflected mode of production.

The fact that social consciousness separates from immediate activity and that its forms acquire a relatively independent existence, as well as its compulsory, binding character in relation to the members of society, was recorded long before Marx. It underlies not only the religious worldview but all of idealist philosophy.

The becoming of the idea of communism required a critique both of this fact itself and of how it is interpreted in theoretical thinking. The critique of social consciousness, not merely of its own era but of any society in which the exploitation of man by man prevails, played a crucial role in the becoming of the ideological content of communism.

Marxism, as is well known, went through the school of criticism of Hegel's philosophy and of Feuerbach's critique of the Christian religion—a form of social consciousness that was dominant in Europe throughout the Middle Ages and continues to claim dominance in bourgeois society as well. Just how seriously Marx took Feuerbach's critique of religion can be seen from the fact that he speaks of nothing less than Feuerbach's great achievement. What, then, with respect to our question, does this achievement consist of, according to Marx? “(1) The proof that philosophy is nothing else but religion rendered into thought and expounded by thought, i.e., another form and manner of existence of the alienation of the human essence; hence equally to be condemned (this is a very important remark concerning philosophy as a form of social consciousness); (2) The establishment of true materialism and of real science, since Feuerbach also makes the social relation of ‘man to man’ the basic principle of the theory.”

These are the merits of Feuerbach's conception, but it is not entirely consistent, and Marx specifically pointed this out: “Feuerbach starts out from the fact of religious self-estrangement, of the duplication of the world into a religious, imaginary world and a real one. His work consists in resolving the religious world into its secular basis. He overlooks the fact that after completing this work, the chief thing still remains to be done. For the fact that the secular basis detaches itself from itself and establishes itself in the clouds as an independent realm can only be explained by the inner strife and intrinsic contradictoriness of this secular basis. The latter must itself, therefore, first be understood in its contradiction and then, by the removal of the contradiction, revolutionised in practice. Thus, for instance, once the earthly family is discovered to be the secret of the holy family, the former must then itself be criticised in theory and transformed in practice by revolutionary means.”

Man’s earthly life—material production—thus becomes for Marx the subject worth studying in order to grasp the essence of the earthly roots of consciousness. “Religion in itself is without content, it owes its being not to heaven but to the earth, and with the abolition of distorted reality, of which it is the theory, it will collapse of itself.” (Marx to Arnold Ruge in Dresden, 30 November 1842)

The same applies to the other forms of social consciousness: the already mentioned philosophy, as well as morality, law, and even science and art, must be understood as aspects of human practice—from practice itself, from its character and content, rather than from themselves. Marxism draws a clear distinction between the basis of social relations—the unity of the productive forces and the relations of production—and the superstructure derived from them, to which social consciousness belongs.

For example, in his preparatory manuscripts for Capital, Marx writes the following about law: “If these trivialities (the representation of law held by the economists contemporary to Marx—ed.) are reduced to their real content, they say more than their preachers know. Namely, that each form of production creates its own legal relations, form of government, etc. The crudity and lack of comprehension lies precisely in that organically coherent factors are brought into an accidental relation, into a merely reflective connection. All the bourgeois economists are aware of is that production can be carried on better under the modern police than, e.g., under club law. They forget only that club law is also law, and that the right of the strongest persists in another form even in their ‘rule-of-law state’.”

The “earthly roots” of social consciousness must not only be understood but also changed—the character of social production must be changed. And the fact that forms of social consciousness grow out of the mode of social production shows that matters with them are by no means as simple as they might seem at first glance. For example, law can be fully sublated, and the need for it as an external regulator overcome, only under definite social conditions—under communism developed on its own foundation. The same applies to overcoming religion. It cannot simply be abolished by declaring it false. It must be sublated, preserving all the achievements of human culture that developed in religious form, and done away with through a radical change in the character of the human activity carried out day after day.

Of course, in a small pamphlet on social consciousness under communism, we can speak only in general terms. But even so, one crucial circumstance must be noted concerning the class-divided society that precedes communism. Just as people occupy different positions in the production process according to their class affiliation, so too does social consciousness have a class character. This does not mean that there is no common field of consciousness, or no objective truth, goodness and beauty that do not depend solely and exclusively on the class position of their bearer. It means that this very field of consciousness is a field of class struggle. It means that the economically ruling class also rules in consciousness. That the consciousness of the ruling class is the ruling consciousness in class society. That is why Lenin, following Marx and Engels, speaks of the importance of struggle in the field of consciousness—of theoretical struggle as a special form of class struggle.

In his pamphlet “What Is to Be Done?”, Lenin devoted an entire section to the significance of the proletariat’s theoretical struggle and to Engels’s view on this question. The following remark by Lenin strikes us as highly relevant today: “Anyone who is at all acquainted with the actual state of our movement cannot but see that the wide spread of Marxism was accompanied by a certain lowering of the theoretical level. Quite a number of people with very little, and even a total lack of, theoretical training joined the movement because of its practical significance and its practical successes. We can judge, therefore, how tactless ‘Rabocheye Dyelo’ is when, with an air of triumph, it quotes Marx’s statement: ‘every step of real movement is more important than a dozen programmes.’ To repeat these words in a period of theoretical disorder is like shouting ‘many happy returns of the day!’ at the sight of a funeral procession. Moreover, these words of Marx are taken from his letter on the Gotha Programme, in which he sharply condemns the eclecticism displayed in the formulation of principles: if you must unite, Marx wrote to the party leaders, then enter into agreements to satisfy the practical aims of the movement, but do not allow any bargaining over principles, do not make theoretical ‘concessions.’ This was Marx’s idea, and yet there are people among us who, in his name, seek to belittle the significance of theory! Without revolutionary theory there can be no revolutionary movement. This idea cannot be insisted upon too strongly at a time when the fashionable preaching of opportunism goes hand in hand with an infatuation with the narrowest forms of practical activity.”

Turning to Engels’s legacy, Lenin cites a lengthy passage, which we consider it necessary to reproduce here as well: “The German workers have two important advantages over those of the rest of Europe. First, they belong to the most theoretical people of Europe, and they have retained that sense of theory which the so-called ‘educated’ classes of Germany have almost completely lost. Without German philosophy, which preceded it, particularly that of Hegel, German scientific socialism—the only scientific socialism that has ever existed—would never have come into being. Without a sense of theory among the workers, this scientific socialism would never have entered their flesh and blood as much as is the case today. What an immeasurable advantage this is may be seen, on the one hand, from the indifference towards all theory which is one of the main reasons why the English working-class movement crawls along so slowly in spite of the splendid organisation of the individual trades; and on the other hand, from the mischief and confusion wrought by Proudhonism, in its original form among the French and the Belgians, and in the form further caricatured by Bakunin among the Spaniards and the Italians.

The second advantage is that, chronologically speaking, the Germans were about the last to come into the workers’ movement. Just as German theoretical socialism will never forget that it rests on the shoulders of Saint-Simon, Fourier and Owen—three men who, for all their fantastic notions and all their utopianism, stand among the most eminent thinkers of all time, and whose genius anticipated countless truths whose correctness we are now proving scientifically—so the German practical workers’ movement must never forget that it has developed on the shoulders of the English and French movements, that it was able simply to turn to its own account their dearly bought experience and to avoid now the mistakes that, at the time, they for the most part could not avoid. Where would we be now without the example of the English trade unions and the French workers’ political struggles, without the colossal impetus given above all by the Paris Commune?

It must be said to the credit of the German workers that they have exploited the advantages of their situation with rare understanding. For the first time since a workers’ movement has existed, the struggle is being conducted pursuant to its three sides—the theoretical, the political and the practical-economic (resistance to the capitalists)—in harmony and in its interconnections, and in a systematic way.”

This passage shows Engels’s attitude towards theoretical struggle. But class struggle in the sphere of social consciousness is not exhausted by theoretical struggle alone, just as social consciousness is not confined to conceptual thinking. The struggle is also waged in the sphere of aesthetic activity, in the sphere of morality, and with respect to law and religion. This, too, must not be underestimated. Nevertheless, the theoretical work by which the proletariat acquires and develops its own class consciousness is the key to all other forms of consciousness, since it makes it possible to understand them.

It is also important to note that in class society, within social consciousness, the universally human is realised through the class-specific, whether or not people understand this. The consciousness of the revolutionary class, which at a given historical stage expresses the universally human interest, is both the driving force of history and the means by which this universally human content is realised. In this sense, not only morality, religion, philosophy and law, but also science and art are ideology as well, since both the aesthetic assimilation of the world in art and cognition in science are conditioned, and therefore limited (the freedom present in them is limited), by the historical movement of society and the balance of forces in the class struggle.

In this connection, I would like to turn to how Marx and Lenin resolved the problem of the subject of thinking. It should be recalled that Marx divided political economy into classical and vulgar, using a quite definite criterion to do so. He points out that classical political economy, unlike vulgar political economy, was a stage in the development of an understanding of its subject matter. He clearly defines the partisanship, that is, the class character, of this science. Both classical and vulgar political economy are bourgeois sciences, serving the bourgeoisie and expressing in theory the direction of its development. As soon as this development ceases to coincide with the development of society and ceases to embody that development—in other words, as soon as the bourgeoisie ceases to be a revolutionary class—political economy becomes vulgar. This means that it is no longer developing thought, but, on the contrary, becomes an obstacle to the development of thought.

Of course, the issue here is not the personal qualities of individual political economists. Marx convincingly shows that the subject of thinking is neither individual theorists nor society as a whole, but precisely a class. Moreover, a class is the subject of thinking only for as long as it is the subject of historical action.

We find the same idea in Lenin when he speaks of parties in philosophy, of which there can be only two: the party that thinks and the party that hinders thinking. And this is precisely because one party is the party of the class that is the subject of historical action, and therefore also the subject of thinking, while the other is not. This applies not only to philosophy but to ALL sciences without exception, including the natural sciences. Although it is less obvious there, it becomes so difficult, for example, to obtain funding for research that promises no returns even in the long run that such cases can quite reasonably be seen as episodes of a veiled class struggle. This remains relevant for as long as society develops through the struggle of classes. It follows that objective truth is always tendentious and has a class character. The same applies to the objectivity of goodness and beauty.

In September (August, Old Style) 1917, Lenin expressed a well-known and, in our view, very important idea in the spirit of what has been written above, one that shows how the proletariat relates to its class consciousness in revolutionary epochs, given that the aim of the proletariat, unlike that of other classes, is not at all to conquer and consolidate its own class domination, but, having conquered that domination, to abolish all class domination and every division of society into classes altogether: “Let us be unswerving,” says Lenin, “in submitting the slightest doubts to the judgement of the class-conscious workers, to the judgement of our Party. We trust it, we see in it the mind, honour and conscience of our epoch, and in the international alliance of revolutionary internationalists we see the sole guarantee of the liberation movement of the working class.” “The domination of a class is impossible without its domination in the sphere of consciousness,” Lenin emphasises. In order to win, the revolutionary class must, in the sphere of consciousness, measure itself not by someone else’s yardstick but by its own, one determined by its historical mission. Having wrested domination in this sphere from the bourgeoisie and armed itself with the most advanced achievements of its own class consciousness, the proletariat becomes a genuine material force of historical transformation: “The weapon of criticism cannot, of course, replace criticism by weapons, material force must be overthrown by material force; but theory also becomes a material force as soon as it has gripped the masses,” Marx wrote in “Contribution to the Critique of Hegel’s Philosophy of Law”.

So far, however, we have been speaking not of social consciousness under communism, but of the state of affairs from which we must start in the transition to it. What, then, does communism as a science insist upon with regard to social consciousness?

Communism insists on overcoming the social division of labour and the alienation of human beings from other people and from the process and the result of their own activity. And this entails corresponding changes in the sphere of social consciousness as well. The alienation of the forms of social consciousness, both from one another and directly from human activity, will also be overcome.

The ideal active capacities of human beings, which are “locked up” within the forms of social consciousness as spheres of aesthetic, moral and intellectual activity separated from material production, will become the possession of each and every one.

The question of social consciousness under communism is a question of practice. It is a question of planning, and of carrying out what has been planned in material activity, in such a way that corresponding changes also take place in the spiritual sphere. By changing their activity, people also change their consciousness as the ideal side of that activity. But under communism, the very change in practice is conscious in character. That is why here it is no longer simply that social being determines social consciousness; rather, social consciousness, as the consciousness of a free association of people acting in concert, determines the being that determines this consciousness.

The socialisation of all spiritual culture, and moreover in accordance with the logic of that culture, becomes a most acute necessity for the transition to communism. That is why, in the theory of communism too, the problems of social consciousness make sense only if they are understood as questions of the becoming of communism.

In communist activity, unlike in alienated labour, the ideal as that which corresponds to an ideal and the ideal as that which pertains to the “world of ideas” coincide empirically in the last instance (precisely in the last instance). The ideal thus appears at once as the reflection (representation, expression) of the material in the material—an active reflection of objective reality—as the ideal goal of activity, to which that activity must conform, and as an aesthetic ideal. Herein lies the most important moment in overcoming the division of labour and its alienation in every sense of the word, and the intrinsic value of activity is affirmed regardless of what exactly is being done.

It should be noted here that it would be wrong to claim that the conscious and material moments of human activity are on an equal footing in labour as such, regardless of its historical specificity, let alone at the level of the activity of society as a whole. Of course, the ideal moment is present in any activity as the moment of goal-setting, unless goal-setting lies outside that activity (for a worker who merely sells his labour power, his own goal does not coincide with the goal of production at all), but that is not the point. In any case, the mode of production of material life determines the mode of production of the ideal. And this has to be reckoned with.

This matters in practical terms because, although an idea does become a material force once it has gripped the masses, the idea itself is the result of the conscious grasp of material forces in the course of their action (the result and an aspect of the interaction of the material with the material). In other words, to change the mode of production of ideas, one must change the mode of production of material life. Only this way, and not the other way round. The other way round, nothing will come of it.

For example, however much we propagandise dialectics today (that is, act ideally), it will not thereby become everyday thinking, and the production of dialectical thinking will for now remain a cottage industry, so to speak. This cottage industry is, of course, very important, and much of the future depends on it. But we must understand that dialectically thinking people will not be a mass product under the present mode of production of material life and the mode of production of consciousness that corresponds to it.

For the problem of communism, it is not enough simply to establish the direction in which man’s material life passes into consciousness (what is primary and what is secondary); one must also put forward ideas for practically accomplishing this transition. And here particular ideas on this score are very important. Ideas for shaping various empirical forms of communism are important, since the material never passes into consciousness all at once and in its entirety, but always does so in certain particular forms. Ideas such as Makarenko’s, for example, are important: love is an ordinary matter that must be properly organised.

However enormous the role played by the ideal moment of activity, what is always at issue is the life and activity of material man in the material world. The ideal is always an attribute of the material activity of social man and corresponds to its movement. The “realm of freedom”, total communism, is characterised, among other things, by the conscious production of social consciousness on a scale that is likewise social.

Isolated breakthroughs in the sphere of social consciousness that go beyond the bounds of capitalism are already taking place. They were discussed (among other things) in the second chapter of this work. This is the production of the ideal in the creation of modern technologies. But that is not the whole story: not all of the ideal belongs to social consciousness. This is no longer simply ideal activity in its alienation, which is superstructural in character, nor is it the ideal confined within the limits of professional cretinism. It is the ideal moment of the production of human life in general. And precisely how this life is produced determines (as with databases and free software, one and the same process is at work here) how a new, immediate, unalienated social consciousness is created, even if only in local areas for now. This is a consciousness produced not only, and not so much, by the mere fact of producing a program as by the very character of its production and functioning.

Social consciousness can determine social consciousness only by determining social being. Moreover, the latter must not be understood as something separate from the whole human being, as something ideal in general, but only as a moment of the activity of the material human being in the equally material world. The point lies precisely in the conscious production of the whole human being: one who acts, feels and thinks universally. Such activity, regardless of what exactly is being done, will appear as the realisation of communism in a directly collective form and will be activity aimed at producing freedom, goal-setting, sensuousness, morality and thinking.

By and large, this presupposes that any activity under communism must, among other things, be the production of ideas and of the ideal of future life, the production of the human being.
\ifSubfilesClassLoaded{\setcounter{page}{62}}{}
\section{Conclusion}
\label{sec:11_conclusion}

Only now, having outlined in broad terms the main elements of communism and shown the seeds from which it can grow into a social system, do we consider ourselves entitled to move on to definitions. Otherwise they would be dead, meaningless abstractions rather than an answer to the question posed in the title: “What is communism?”

Communism (communist society) is a free association of people who rationally create their own lives and consciously posit their own free activity.

Communism is the directly collective form of free activity, having become total as the form of a society that creates its own development.

Communism is, as Marx and Engels put it in the “Manifesto”, “an association, in which the free development of each is the condition for the free development of all”.

The time has come to repeat once again a slogan that has not aged one bit:

WORKERS OF ALL COUNTRIES, UNITE!

\end{document}
